% The part on the methods behind Scholar's defaults (#310, #311).
% \input from scholar.tex after the Tutorials part.  The chapters are
% plain LaTeX, not literate programs: they hold no code.  The sources
% they cite are in doc/bibliography.bib, each with its provenance block.

The tutorials show how to run Scholar; this part says how to search and
screen with it, and on what grounds.
\Cref{ch:procedure} is the procedure, one step after another, each with
the command that carries it out and a pointer to the chapter that backs
it.
\Cref{ch:measured} reports what we measured on searches and screens of
our own: what a limit on a search lost, what screening everything added,
how the classifier behaved in batches of different sizes, and how the
tiers of the literature search behind the later chapters turned out.
\Cref{ch:lit-search} describes that literature search once, and the
chapters after it each take one claim the procedure rests on and give
what the literature says for and against it, at the strength the sources
carry.

Two kinds of evidence meet here and they are kept apart.
Our measurements are a few cases from one author's work, screened with
one model or one agent; they show what happened in those cases and are
reported as such.
The literature chapters cite only sources that were read, each with a
verbatim quote that the bibliography's source file records beside the
entry, together with how the source was found, why it was picked and
whether it was read in full or only in its abstract.

\chapter{How to search and screen}
\label{ch:procedure}

This is the procedure in the order it is carried out.
Each step names the command that does it and the chapter or section
that gives its grounds: a measurement of our own (\cref{ch:measured}),
the literature (\cref{ch:stopping,ch:abbreviated,ch:databases,%
ch:llm-screening,ch:reporting}), or the code.
The commands are those of this version of Scholar; \texttt{CLAIM} stands
for a session's name and \texttt{\dots} for what the reader fills in.

\begingroup\raggedright
\begin{description}
\item[Check the providers]
  \texttt{scholar providers check} probes each configured database with
  a one-result request and reports it as ok, rate limited, key rejected
  and so on (\cref{sec:provider-health}).
  A database that fails is not searched, and the method says so; a
  search that fails later is recorded as failed, never as zero hits.
  In health research no single database held what reviews included, and
  outside health the right databases differ (\cref{ch:databases}).

\item[One session per claim]
  Every query for one claim goes into one session,
  \texttt{scholar search \dots\ -n CLAIM}, so a work that several queries
  return is screened once.
  In a book chapter screened with one session per query, over a third of
  the screening was of works already screened in another session of the
  same claim (\cref{sec:measured-chapter}).

\item[Several wordings, fixed beforehand]
  Each concept of the claim gets several wordings (synonyms, the noun and
  the verb, the bare word beside the phrase), and each claim gets counter
  queries that reuse every concept term of the support queries.
  The wordings and the studies the queries must return
  (below) are written down before the first search.
  In the literature search of \cref{ch:lit-search}, the check by meaning
  found works the first wordings had missed, and new wordings of the
  same concepts brought several of them back.

\item[Every wording on every live provider, without a limit]
  \texttt{scholar search \textquotesingle{}\dots\textquotesingle{} -n CLAIM -{}-in title-abstract -p openalex
  -p s2 -p scopus -p wos -p ieee -p arxiv}, with no \texttt{-{}-limit}, so
  every result is fetched (\cref{sec:no-limit-default}).
  With \texttt{-{}-in title-abstract}, OpenAlex matches title and abstract
  instead of its default fields, which include the full text, and Scopus
  gets the query as \texttt{TITLE-ABS-KEY(\dots)}; Web of Science wraps a
  plain query as \texttt{TS=(\dots)}, its title, abstract and keyword
  fields (\cref{sec:providers-search-in}).
  A limit of 40 per query and database hid four of the sources a review
  later cited, one of them the only review against a hypothesis
  (\cref{sec:measured-limit}); what the literature says about shortened
  searches is in \cref{ch:abbreviated}, and the reporting guidelines ask
  that any limit be stated and justified (\cref{ch:reporting}).

\item[Check by meaning, then repair]
  \texttt{scholar sessions check CLAIM -s \textquotesingle{}<the claim as a question>\textquotesingle{}}
  looks each work of an OpenAlex semantic search up in the sessions
  (\cref{sec:sessions-check}).
  A work that bears on the claim and is in no session is a gap in the
  wordings: a new query by concept is added and run on every provider.
  In the literature search, 29 of the 50 works the check returned for one
  question were in no session, and two repair queries brought back 17 of
  them; of the works the checks returned that were in the sessions, 48 to
  73\,\% bore on the question, against 5 to 15\,\% of the first 100
  places of the keyword answers (\cref{sec:measured-tiers}).
  The check is a ranked list of at most 50 works, so it never counts
  anything and never shows that something does not exist.

\item[The named studies must come back]
  \texttt{scholar sessions check CLAIM -s \textquotesingle{}\dots\textquotesingle{} -k DOI} also locates
  each study named beforehand in the sessions.
  A query for a concept that returns none of the known studies on that
  concept is reworded and run again.

\item[Screen in tiers]
  The first tier is screened: the first 100 places of each query in each
  answer ordered by relevance (\texttt{sessions show} prints the order a
  search recorded), what the check by meaning returned, and the named
  studies.
  \texttt{scholar sessions export CLAIM -f screen-sheet -{}-pending
  -{}-top 100 -o BASE} writes the first part as a sheet and its counts to
  \texttt{BASE-tier.json} (\cref{sec:sheet-tier}); the verdicts go back
  with \texttt{scholar sessions decide CLAIM -{}-source agent
  -{}-reasons-file F} or are made in the review screens.
  A tier decides what is read first, not what exists: stopping a ranked
  screen at a fixed place or by a simple rule often misses relevant
  records (\cref{ch:stopping}), and relevant records sit deep in ranked
  lists (\cref{ch:databases}).

\item[Sample the rest]
  \texttt{scholar sessions export CLAIM -f screen-sheet -{}-pending
  -{}-top 100 -{}-rest-sample 60 -{}-seed S -o BASE} draws a seeded random
  sample of the pending records the tier leaves out and records the draw
  in the session (\cref{sec:rest-sample-sheet}); once it is screened,
  \texttt{scholar sessions tier-sample CLAIM -{}-report} gives the share of
  bearing records with its 95\,\% modified Wilson interval
  (\cref{sec:sessions-tier-sample,sec:llm-sample-interval}).
  That share is the one statement the method makes about what was not
  screened (\cref{sec:stopping-sample}).

\item[The classifier, and what its run line says]
  \texttt{scholar llm context CLAIM \textquotesingle{}\dots\textquotesingle{}} states the question and the
  criteria; \texttt{scholar llm classify CLAIM -{}-top 100 -n N} classifies
  the first tier's pending papers (\cref{sec:llm-classify-tier}).
  The run ends with one line of the form \enquote{Requests: 1 batch + 1
  retry of 2 papers; left pending: 0; tokens: \dots}: how many
  papers had to be asked again, how many were left without a usable
  answer, and the tokens the model reported (\cref{sec:classifier-run-cost}).
  In our test on 50 papers, runs with batches of different sizes agreed
  on 45 to 48 of the 50 kept-or-discarded verdicts, two runs with the
  same batch size on 49 (\cref{sec:measured-classifier}); what is known
  about screening with a
  language model is in \cref{ch:llm-screening}.

\item[A second screen where the result matters]
  \texttt{scholar llm classify CLAIM -{}-second-screen -{}-over machine}, or
  another screener's verdicts with \texttt{scholar sessions screen},
  stores a second verdict beside each decision without changing it
  (\cref{sec:llm-classify-second-screen,sec:sessions-screen});
  \texttt{scholar sessions agreement CLAIM} and \texttt{scholar sessions
  conflicts CLAIM -{}-unresolved} give the agreement and the papers to
  settle (\cref{sec:sessions-agreement,sec:sessions-conflicts}).
  One screener misses more relevant studies than two
  (\cref{sec:abbreviated-single}), and two runs of one model on the same
  papers did not give the same verdicts (\cref{ch:llm-screening}).

\item[Write what was done and what was not]
  The method states what was searched (databases, queries, dates, hits
  per query and database), what was screened (the tier's definition and
  size), what was not (how many records, and the sample's size, seed,
  count, rate and interval), and who screened.
  \texttt{scholar sessions show CLAIM} gives the hits per query and
  database, and \texttt{scholar sessions export CLAIM -f table
  -{}-bearing-only -{}-cited-from refs.bib} the table of records that bear on
  the claim (\cref{sec:audit-table,sec:audit-cited-from}).
  A negative result is worded to match: \enquote{none among the $N$
  screened of $M$ returned; a sample of $n$ of the rest held none
  (interval)}, never \enquote{none found}.
  The reporting guidelines behind this are in \cref{ch:reporting}.
\end{description}
\endgroup

\chapter{What was measured}
\label{ch:measured}

% Author's note: every number in this chapter is from the analysis notes
% of 2026-11-23 (campaign-2026-11d/cap-analysis.md, scripts listed in
% campaign-2026-11d/measurements/README.md), except the tiers of the
% literature search in sec:measured-tiers, which are from that search's
% report (campaign-2026-11d/lit-screening/work/report.md).  The scripts
% read private sessions and are not in this repository.

The defaults of \cref{ch:procedure} were set after four looks at
searches and screens of our own, all made on 2026-11-23.
Two are cases from one author's work: a review of one security question,
searched first with a limit and later without one
(\cref{sec:measured-limit}), and one book chapter's sixteen sessions,
screened in full (\cref{sec:measured-chapter}).
The third is a test of the classifier on fifty papers
(\cref{sec:measured-classifier}), the fourth the screening tiers of the
literature search of \cref{ch:lit-search} (\cref{sec:measured-tiers}).

What these measurements are matters for how far they carry.
They are a few cases, from one author's searches, with one language
model and one agent screener; they show what happened in these cases,
not what happens in general.
Where they say a work bears on a claim, that is the verdict of a
classifier or an agent, not of a person who checked it; where they say a
work was cited, that is what the authors chose from what they had found.
The scripts that computed them read the author's own sessions, which are
not published, so the numbers cannot be recomputed from this repository;
the literature search is the exception in part, since its search record
and the records it judged bearing are printed in the chapters that use
it.

\section{A limit on the search, and no limit}
\label{sec:measured-limit}

A review of one security question was searched with a limit of 40
results per query and database, written up, and later searched again
with the same queries and no limit.
We joined the two sets of exported records by paper.

\begin{description}
\item[Records] The search without a limit returned 8,267 records
  against 1,873 with it, 6,394 more.
\item[Bearing] Of the new records, 139 (2.2\,\%) were judged to bear on
  the question, against 7 to 12\,\% of the records the limited search
  had returned (the share differs by sub-question).
\item[By database] The new records came from OpenAlex (4,440, of which
  2.4\,\% bearing), IEEE Xplore (1,036; 0.4\,\%), Scopus (411; 5.4\,\%),
  arXiv (306; 0.3\,\%) and Web of Science (201; 3.5\,\%).
\item[Cited sources] The bibliography gained seven sources.
  Four were not among the limited search's records, two were (one of
  them discarded there), and one is in neither export.
  Of the cited sources found in the exports, 52 of 56 were already among
  the limited search's records.
\item[Conclusions] No result was overturned.
  One of the four sources the limit had hidden was the only review
  against one of the hypotheses.
  With the limit, the text had said that the search found no review
  against that hypothesis, which was true only of the limited search.
\end{description}

Two further changes to the review came from elsewhere: a study that no
query returned, found by comparing with another report, and four
sentences narrowed and two sources dropped after reading the cited works
in full.

\section{One book chapter's sixteen sessions}
\label{sec:measured-chapter}

A book chapter backs three claims with sixteen sessions, 36,808 records
in all, every one of them screened.
Before that screen the chapter had been written from a keyword pass,
without screening the records, followed by the check by meaning and the
repair queries it led to.

\paragraph{What the full screen added.}
The keyword pass gave 15 sources and the repair queries 4 counter
sources.
The finished chapter cites 16 works: 12 of the 15 and the 4 from the
repair queries.
The screen of all 36,808 records added no cited work.
The conclusion for each claim stayed as it was (one claim backed, one
narrower than first stated, one not backed); four works moved from
supporting to qualifying and one phrase was dropped, all after reading
the works in full.

\paragraph{Repeated screening.}
The 36,808 records were 20,678 distinct works.
Within one claim, 37\,\% of the screening was of a work already in
another session of the same claim (42\,\% and 34\,\% for the two claims
with several sessions), because each query had a session of its own.

\paragraph{A limit, simulated per list.}
From the cached answers of twelve queries, 68 lists in all, we took the
first $L$ of every list and counted what was left
(\cref{tab:measured-cap}).
Semantic Scholar's lists were ordered by citations, as its bulk search
orders a limited search.
The one cited work beyond place 500 was at place 512 of an OpenAlex
list.
Of the 16 cited works, 14 stood within the first 40 places of some
list, 13 of them in an OpenAlex or arXiv list and one in a Web of
Science list.

\begin{table}
  \begin{sidecaption}{%
      What the first $L$ places of each list would have held, of 19,607
      works in 68 lists of one book chapter, simulated from the cached
      answers.
      Works: the share of all works.
      Cited: of the chapter's 16 cited works.
      Bearing: the share of the works judged bearing.
    }[tab:measured-cap]
    \begin{tabular}{@{}rrrr@{}}
      \toprule
      $L$ & Works & Cited & Bearing \\
      \midrule
      40 & 1,551 (8\,\%) & 14 & 24\,\% \\
      100 & 16\,\% & 15 & 42\,\% \\
      200 & 27\,\% & 15 & 58\,\% \\
      500 & 48\,\% & 15 & 81\,\% \\
      1,000 & 66\,\% & 16 & 92\,\% \\
      \bottomrule
    \end{tabular}
  \end{sidecaption}
\end{table}

\paragraph{Per list or per question.}
For the two claims with several queries we took the lists of the
databases that order by relevance (OpenAlex, arXiv and IEEE Xplore;
14 lists from 5 queries and 19 from 7) and compared two ways to cut them.
The first $K$ places of every list, united, held 481 works and 8 of the
16 cited at $K=20$, 897 works and 14 cited at $K=40$, and 1,991 works
and 15 cited at $K=100$.
A fixed number $N$ of works per claim, taken in order of each work's best
place in any list, held 4 cited at $N=100$ (200 works for the two claims;
the 100th work had best place 10 and 7), 7 at 200, 10 at 300, 14 at
500 (1,000 works) and 15 at 1,000 (2,000 works).
So in this case a number per list reached as many cited works with
slightly fewer records than a number per claim (14 with 897 works against
1,000), and a small number per claim stopped near the top of every list:
100 per claim reached only the tenth place.

\paragraph{Where the bearing works sit.}
\Cref{tab:measured-density} gives the share of bearing works by place
in each database's lists.
In OpenAlex, arXiv and IEEE Xplore the share falls with the place; in
Scopus and Web of Science it does not, and those lists were ordered by
date: every neighbouring pair of Scopus records was in non-increasing
year order (mean year 2026.0 in the first quarter of a list, 2022.0 in
the last), and 89\,\% of Web of Science's pairs were, while OpenAlex,
arXiv and IEEE Xplore were not ordered by year.
Scholar sent no sort order to Scopus or Web of Science then; it now asks
both for relevance (\cref{sec:answer-orders}).

\begin{table}
  \begin{sidecaption}{%
      Share of works judged bearing, by place in each database's lists,
      in one book chapter's sessions.
      Semantic Scholar's lists were ordered by citations, Scopus's and
      Web of Science's by date.
    }[tab:measured-density]
    \begin{tabular}{@{}lrrrrr@{}}
      \toprule
      & 1--40 & 41--100 & 101--200 & 201--500 & 501-- \\
      \midrule
      OpenAlex & 42\,\% & 33\,\% & 29\,\% & 22\,\% & 11\,\% \\
      arXiv & 38\,\% & 24\,\% & 14\,\% & 10\,\% & 3\,\% \\
      IEEE Xplore & 23\,\% & 10\,\% & 8\,\% & 6\,\% & 2\,\% \\
      Semantic Scholar & 30\,\% & 26\,\% & 25\,\% & 23\,\% & 14\,\% \\
      Scopus & 26\,\% & 29\,\% & 29\,\% & 27\,\% & 18\,\% \\
      Web of Science & 24\,\% & 28\,\% & 28\,\% & 26\,\% & 34\,\% \\
      \bottomrule
    \end{tabular}
  \end{sidecaption}
\end{table}

\paragraph{What each database contributed.}
\Cref{tab:measured-providers} counts, per database, the works it
returned in the sixteen sessions, those no other database returned, the
bearing and the cited ones.
OpenAlex and arXiv together returned all 16 cited works.
The share of bearing works rose with the number of databases that
returned a work: 4\,\% for one, 8\,\% for two, 10\,\% for three and
19\,\% for four.

\begin{table}
  \begin{sidecaption}{%
      What each database contributed to one book chapter's sixteen
      sessions: works returned, works only it returned, works judged
      bearing, cited works, and cited works only it returned.
    }[tab:measured-providers]
    \begin{tabular}{@{}lrrrrr@{}}
      \toprule
      & Works & Only it & Bearing & Cited & Only it, cited \\
      \midrule
      OpenAlex & 13,450 & 6,345 & 1,146 & 14 & 4 \\
      Semantic Scholar & 10,191 & 2,858 & 850 & 10 & 0 \\
      arXiv & 5,019 & 1,757 & 259 & 8 & 0 \\
      Scopus & 3,537 & 496 & 540 & 5 & 0 \\
      Web of Science & 2,252 & 279 & 404 & 4 & 0 \\
      IEEE Xplore & 1,436 & 612 & 60 & 0 & 0 \\
      \bottomrule
    \end{tabular}
  \end{sidecaption}
\end{table}

\paragraph{Cost.}
A pilot classification in one of the sessions used about 158,000 tokens
of a Claude Sonnet model per 100 records.

\section{The classifier in batches}
\label{sec:measured-classifier}

We gave \texttt{scholar llm classify} the same 50 papers of one claim, all
with abstracts and no examples in the prompt, in three shapes: one batch
of 50, five batches of 10 (twice), and one paper per request.
The model was \texttt{anthropic/\allowbreak claude-sonnet-5-5}, and every request
was a real call to it (\cref{tab:measured-batches}).

\begin{table}
  \begin{sidecaption}{%
      One classifier, 50 papers, three batch sizes.
      Retried: requests whose answer had to be asked again.
      Answered: papers with a usable answer at the end.
      Seconds: the run's duration.
      Conf.: the mean confidence the model gave.
      Characters: of the prompts sent, retries included.
    }[tab:measured-batches]
    \footnotesize
    \begin{tabular}{@{}lrrrrrrr@{}}
      \toprule
      Batches & Requests & Retried & Seconds & Answered & Kept &
      Conf. & Characters \\
      \midrule
      1 of 50 & 1 & 0 & 57 & 50 & 26 & 0.74 & 104k \\
      5 of 10 & 5 & 1 & 91 & 50 & 27 & 0.78 & 131k \\
      5 of 10, again & 5 & 1 & 85 & 50 & 26 & 0.78 & 131k \\
      50 of 1 & 50 & 13 & 513 & 50 & 26 & 0.79 & 328k \\
      \bottomrule
    \end{tabular}
  \end{sidecaption}
\end{table}

No paper was skipped in any shape.
Every retry was for a first answer without a category tag, and after the
one retry every paper had a usable answer.
The two identical runs in batches of ten agreed on 49 of the 50
kept-or-discarded verdicts and on 46 of the categories; runs with
different batch sizes agreed on 45 to 48 verdicts and 38 to 43
categories.
Six papers changed status in at least one run, in no one direction: 26
or 27 papers were kept in every run.
Against the verdicts stored before (an earlier classifier run in batches
of 100 for 46 papers, an agent for 4; not a reference standard) the runs
agreed on 44 to 47 statuses and 38 to 43 categories.

\paragraph{A prompt with the papers last, withdrawn.}
A version of the prompt that put the papers after the instructions was
tried on the same papers.
In one batch of 50 the first answer had no category tag on any of the
50 papers, and each paper was asked again on its own: 138,101 tokens in
and 15,014 out, 153 seconds.
In five batches of 10, three needed no retry (7,200 to 9,500 tokens in)
and two lacked the category on 9 and on 10 of their papers (25,700 and
28,000 tokens in).
After the retries all 50 answers were usable, and they agreed with the
earlier runs on 43 to 46 statuses and 40 to 44 categories.
The prompt went back to its earlier order, and the retry became one
request holding only the papers whose answer was unusable
(\cref{sec:retry-one-paper}).

\paragraph{The revised code.}
With the earlier order and the new retry, one batch of 50 took one
request and no retry: 34,898 tokens in, 6,830 out, 51 seconds, all 50
answered, 28 kept.
Five batches of 10 took five requests and no retry: 39,410 tokens in and
5,845 out in all (7,200 to 9,400 in each), 12 to 14 seconds each, all 50
answered, 28 kept.
The revised runs agreed with the first test's batch of 50 on 48 of 50
statuses and 43 categories, with its batches of ten on 49 and 47, and
with each other on 48 and 43.
Counting all runs: with the earlier prompt order a retry was needed in
none of the batches of 50 and in 2 of 10 batches of 10 in the first test,
and in none of the 6 runs of the revised code; with the papers last, in
3 of 6 runs, each time for most or all papers of the batch.
A paper cost about 835 tokens, in and out together, in a batch of 50, and
about 905 in batches of 10.

\paragraph{Limits.}
Fifty papers, one claim, one model, and one repeat of one shape; the
first test did not record tokens or price.

\section{The tiers of the literature search}
\label{sec:measured-tiers}

The literature search of \cref{ch:lit-search} screened a first tier and
sampled the rest, so it measured the tier as it went.
Its four questions are those of \cref{ch:stopping,ch:abbreviated,%
ch:databases,ch:llm-screening}.
Bearing means judged to support or to qualify the question's claim, by
the one agent that screened the search.

\paragraph{The parts of the first tier.}
\Cref{tab:measured-tier-parts} gives each part's records and how many
bore on the question.
The check by meaning's works were the densest part, 48 to 73\,\%
bearing against 5 to 15\,\% for the first 100 places of the keyword
answers.
Records returned by several databases added little: beyond the first
100 places, 2 of the 49 records that three or more databases returned
for the first question bore on it, and a seeded sample of 60 records
returned by three or more databases, for each of the second and third
questions, held none (and one adjacent record each).
The tier was large because the answers were: 26 of the 35 OpenAlex
answers held more than 100 records.

\begin{table}
  \begin{sidecaption}{%
      The parts of the first tier in the literature search: records
      screened, and in parentheses those bearing on the question.
      First 100: places 1 to 100 of each query's OpenAlex and arXiv
      answers.
      Several databases: returned by at least three (first question) or
      five (the others) databases, beyond the first 100.
      Check: the check by meaning's works that were in the sessions.
      Known: named studies not in the other parts (--: none reported).
    }[tab:measured-tier-parts]
    \begin{tabular}{@{}lrrrr@{}}
      \toprule
      & First 100 & Several databases & Check & Known \\
      \midrule
      \Cref{ch:stopping} & 401 (49) & 49 (2) & 38 (22) & 1 \\
      \Cref{ch:abbreviated} & 675 (36) & 1 (0) & 17 (10) & 1 \\
      \Cref{ch:databases} & 996 (54) & 17 (0) & 21 (10) & -- \\
      \Cref{ch:llm-screening} & 645 (98) & 12 (2) & 45 (33) & -- \\
      \bottomrule
    \end{tabular}
  \end{sidecaption}
\end{table}

\paragraph{The samples of the rest.}
For each question, 60 records were drawn at random from the records left
unscreened and screened as the first tier was (\cref{sec:search-tiers}).
Three of the 240 bore on a question, all three on the first; in the first
tier 254 of 2,788 did (\cref{tab:measured-samples}).
Pooled, 2 of the 172 sampled records that one database returned were
bearing, 1 of 37 returned by two, none of 15 returned by three and none
of 16 returned by four or more.
None of the three would have entered the first tier by any of its
rules: none was in an OpenAlex or arXiv answer to the first question's
queries, none was returned by three or more databases, and none was
among the check by meaning's works.
Two were returned only by Web of Science or Scopus, whose answers were
then not ordered by relevance, and the third by both.

\begin{table}
  \begin{sidecaption}{%
      Records bearing on the question in the seeded sample of the
      unscreened records and in the first tier, with the 95\,\% modified
      Wilson interval of the sample's share
      (\cref{sec:llm-sample-interval}).
    }[tab:measured-samples]
    \begin{tabular}{@{}lllr@{}}
      \toprule
      & Sample & Interval & First tier \\
      \midrule
      \Cref{ch:stopping} & 3 of 60 (5\,\%) & 1.4 to 13.7\,\% & 53 of 454 \\
      \Cref{ch:abbreviated} & 0 of 60 & 0 to 6.0\,\% & 38 of 679 \\
      \Cref{ch:databases} & 0 of 60 & 0 to 6.0\,\% & 58 of 1,003 \\
      \Cref{ch:llm-screening} & 0 of 60 & 0 to 6.0\,\% & 105 of 652 \\
      All & 3 of 240 (1.25\,\%) & 0.3 to 3.6\,\% & 254 of 2,788 \\
      \bottomrule
    \end{tabular}
  \end{sidecaption}
\end{table}

\chapter{How the literature was searched}
\label{ch:lit-search}

% Author's note (verification and reproduction):
%   sessions      : screenlit-q1 .. screenlit-q4 (one per question) and
%                   screenlit-q1-top .. screenlit-q4-top (the first 100
%                   places of the OpenAlex and arXiv answers), private
%                   data directory of the review of 2026-11-23
%   queries       : lit-screening/work/queries.tsv, queries-added.tsv
%   provider check: lit-screening/work/providers-check-2.txt
%   checks        : lit-screening/work/checks/q*.json (sessions check)
%   decisions     : lit-screening/work/decisions/*.tsv, dec_q*.py
%   samples       : lit-screening/work/tier2-sample-q*.tsv (seed
%                   20261123+q), sample60-q2/q3.txt (seed 20261123)
%   report        : lit-screening/work/report.md
%   tables        : doc/methods/tables/*.tex, made from copies of the
%                   sessions with this branch's scholar (matrix cells as
%                   hits_label gives them; audit tables by
%                   generate_audit_table with --bearing-only and
%                   --cited-from doc/bibliography.bib)

The four chapters that follow each answer one question with the
literature.
They share one search, made on 2026-11-23 with Scholar by one agent,
which this chapter describes once; each chapter then gives its own
queries, their hits per database, its tiers and its sample, and the
records it judged to bear on its question.

\section{Questions and claims}
\label{sec:search-questions}

Each question was put as the claim that the procedure would rest on if
the claim held, so that the searches could find evidence against it as
well as for it.

\begin{description}
\item[First question (\cref{ch:stopping})] Stopping a prioritised screen
  before the end of the list is reliable enough: recall targets and
  stopping rules keep the share of missed relevant studies small.
\item[Second question (\cref{ch:abbreviated})] An abbreviated search or
  screen (fewer databases, a result limit, a rapid review, single
  screening) loses few included studies and leaves pooled estimates and
  conclusions unchanged.
\item[Third question (\cref{ch:databases})] A few bibliographic databases
  (Scopus, Web of Science, Google Scholar, OpenAlex, Semantic Scholar), or
  the first results of a ranked source, find nearly all the relevant
  records a review needs.
\item[Fourth question (\cref{ch:llm-screening})] A cheaper first pass
  with a language model (title-only, a small model, a model on the
  confident band) misses few relevant studies compared with human
  screening.
\end{description}

\section{Databases and queries}
\label{sec:search-databases}

\texttt{scholar providers check}, run before the searches, found
OpenAlex, Scopus (with the full record view), Web of Science and IEEE
Xplore working, Semantic Scholar and arXiv answering with rate limits
(HTTP 429) at that moment, and DBLP failing: it answered with an HTML page
instead of data.
DBLP was therefore not searched.
Semantic Scholar stayed rate-limited through the day, and its answers to
four queries of the third question are incomplete; arXiv answered every
search.

Each question has one session, and every query of a question ran on the
same six databases with the same keywords and without a limit:
OpenAlex in title and abstract (\texttt{-{}-in title-abstract}), Semantic
Scholar through its bulk search, Scopus as \texttt{TITLE-ABS-KEY(\dots)},
Web of Science as \texttt{TS=(\dots)}, IEEE Xplore and arXiv.
The sessions record all of this per search (no limit, the title and
abstract filter for OpenAlex, the bulk search for Semantic Scholar), and
every search is dated 2026-11-23.
The query tables of the following chapters give the queries verbatim and
the hits per database as \texttt{scholar sessions show} reports them.
A zero from IEEE Xplore or arXiv is an answer, not a failure: both
answered other queries of the same syntax with hits.
PubMed is not among Scholar's databases; the records from health methods
journals, where most of the evidence turned out to be, came through
OpenAlex, Scopus and Web of Science (the database column of the tables
of bearing records shows which).

Each question has support queries, at least one second support wording,
and counter queries that reuse every concept term of the support
queries with words for failure, missed studies or changed conclusions.
Some queries were added after the check by meaning
(\cref{sec:search-check}); the query tables mark them.
In all, the four questions have 35 queries.

Before the searches, a few studies known to bear on each question were
named, by following the references of two reviews (one on stopping
rules in education, one on database combinations) through OpenAlex.
All of them came back except two: one study on the contribution of
databases was returned only by the second question's queries, though it
bears on the third too, and one older study on searching several
databases was returned by no query and is not cited.

Scholar deduplicated each search's records: records with one DOI, or
without one but with the same title and first author and years at most
one apart, were merged, and a preprint was folded into its published
version when that was unambiguous (\cref{sec:same-work-guard,%
sec:version-dedup}).

\section{The check by meaning}
\label{sec:search-check}

For each question, its claim was put to OpenAlex's semantic search as a
plain question, and each of the 50 works returned was looked up in the
question's sessions (\texttt{scholar sessions check}).
A work that bore on the question and was in no session was taken as a
gap in the wordings, and a query by concept was added for it and run on
every database.
For the first question 29 of the 50 works were in no session; two
repair queries brought back 17 of them, and the 12 left outside were
grey literature, items from 2026 and one preprint on stopping.
For the second question 36 were missing; one repair query returned three
of the key comparisons, and the 33 left outside were mostly general
texts on systematic reviews.
For the third, 29 were missing; one repair query returned one of them, and
26 comparisons of bibliographic coverage stayed outside.
For the fourth, 4 were missing and no query was added; among them was a
comparison of screening by title with screening by title and abstract,
which one query's wording should have returned.
A further repair query for the third question was added because two of
the named studies had not come back on its concepts.

\section{Tiers and samples}
\label{sec:search-tiers}

Not every record was screened.
The first tier was:
\begin{enumerate}
\item the first 100 places of each query's OpenAlex and arXiv answers,
  fetched again with a limit of 100 into a companion session per
  question, because the version of Scholar used then sent Scopus and Web
  of Science no sort order and their answers came newest first, or
  mostly so (\cref{sec:measured-chapter});
\item the records returned by at least three databases, for the first
  question, and by at least five for the others;
\item the works of the check by meaning that were in the sessions;
\item the named studies.
\end{enumerate}
The rest of each question's records, pending in its session and not
decided in the companion session, was left unscreened, and 60 of them
were drawn at random (Python's \texttt{random.Random(20261123 + $q$)}
for question $q$, over the sorted paper ids) and screened as the first
tier was.
For the second and third questions, 60 records of the group returned by
three or more databases were also drawn (seed 20261123) and screened.
\Cref{tab:search-tiers} gives the counts; the intervals are the 95\,\%
modified Wilson intervals that Scholar computes
(\cref{sec:llm-sample-interval}).

\begin{table}
  \begin{sidecaption}{%
      The literature search in numbers: records returned (with
      duplicates across queries and databases), distinct records,
      records screened in the first tier, bearing records among the
      screened (supporting plus qualifying), cited works (read in full
      plus abstract only), and records left unscreened.
    }[tab:search-tiers]
    \begin{tabular}{@{}lrrrrrr@{}}
      \toprule
      & Returned & Distinct & First tier & Bearing & Cited & Unscreened \\
      \midrule
      \Cref{ch:stopping} & 2,264 & 931 & 454 & 23 + 30 & 8 + 3 & 477 \\
      \Cref{ch:abbreviated} & 8,627 & 4,123 & 679 & 12 + 26 & 4 + 11 & 3,384 \\
      \Cref{ch:databases} & 113,942 & 61,348 & 1,003 & 22 + 36 & 7 + 4 & 60,268 \\
      \Cref{ch:llm-screening} & 6,818 & 2,574 & 652 & 63 + 44 & 7 + 3 & 1,910 \\
      \bottomrule
    \end{tabular}
  \end{sidecaption}
\end{table}

\section{Screening and reading}
\label{sec:search-screening}

One AI agent screened every record, from the title and, for anything
about methods, the abstract; no classifier was used and no person checked
its verdicts or a sample of them.
Its decisions are recorded in the sessions as an agent's
(\texttt{scholar sessions decide -{}-source agent -{}-reasons-file}), each
with its own reason.
Each kept record has one category: it supports the claim, or it
qualifies or contradicts it.
Each discarded record is either off topic, with a criterion (not a
screening or search method, no abstract, or not a study), or adjacent (the
right field, another question), with the question's own label for what it
lacks (no evidence on stopping, no comparison, no evidence on retrieval,
no measurement of errors).
Since nobody reviewed the agent's verdicts, the tables of bearing records
in the following chapters mark a record as cited when the text cites it,
not because a person kept it.

The cited works were read in full where the publisher's copy, an open
copy or a preprint could be had; 30 of the 51 were read in full (six of
them as preprints of papers whose published versions were not obtained,
or that are preprints only), 21 only in their abstract, and each chapter
says which.
Twenty works were wanted and not obtained; each chapter names those that
would have borne on its question.

\section{Limits of the search}
\label{sec:search-limits}

One screener, an agent, decided everything; how often it erred is not
known.
Most of the evidence found is from health research, with two or three
studies per question from computing and two from education.
Several sources on language models are preprints.
Four Semantic Scholar answers of the third question are incomplete
(rate limits), and its Scopus answer to one query stopped at 4,991 of
18,141 records (HTTP 400); its OpenAlex answer to the same query holds
9,978 of the 26,809 records OpenAlex reported, and the session does not
mark it incomplete.
The first tier of the third question was larger than planned (1,003
records against 900).
Seven bearing records of the third question were decided in its companion
session, two of them records that the question's full session does not
hold; its table of bearing records therefore comes in two parts
(\cref{sec:databases-hits}).

\chapter{Stopping a ranked screen early}
\label{ch:stopping}

% Author's note (verification and reproduction): session screenlit-q1
% (and screenlit-q1-top); queries q1a-q1g, q1x, q1y; check by meaning
% checks/q1*.json; sample tier2-sample-q1.tsv (seed 20261124); see the
% block at the top of search.tex.  Sources: doc/bibliography.bib, keys
% Repke2026 Callaghan2020 Kempny2026 Campos2024 Koenig2024 Boetje2024
% Hamel2020 Gates2020 Callaghan2024 Slade2026 Briscoe2023.

The procedure of \cref{ch:procedure} screens a first tier, samples the
rest and states both; it never stops because the records have grown
thin.
Tools that rank records for screening offer another way: screen in the
order of the ranking and stop when a rule says that enough of the
relevant records have been found.
If such a rule were reliable, the first tier could end where the rule
says, and nothing more would need to be said about the rest.
This chapter asks whether it is: \emph{does stopping a prioritised
screen before the end of the list keep the share of missed relevant
studies small?}
The claim is tested at the strength a procedure would need: a rule that
holds a stated recall target across reviews, not one that sometimes
does.

\section{Method}
\label{sec:stopping-method}

The search is the one described in \cref{ch:lit-search}.
The first question's queries are listed below; \cref{tab:stopping-matrix}
gives their hits per database.
The support queries are q1a to q1e, the counter queries q1x and q1y; q1f
and q1g were added after the check by meaning, with the words that tool
evaluations use (sensitivity, missed records, workload; prioritisation,
filtering, ranking, work saved).
\begin{description}[style=nextline,font=\normalfont\bfseries]
\item[q1a, support query] {\footnotesize\ttfamily\raggedright (\textquotedbl{}stopping rule\textquotedbl{} OR \textquotedbl{}stopping rules\textquotedbl{} OR \textquotedbl{}stopping criterion\textquotedbl{} OR \textquotedbl{}stopping criteria\textquotedbl{} OR \textquotedbl{}stopping method\textquotedbl{}) AND (screening) AND (\textquotedbl{}systematic review\textquotedbl{} OR \textquotedbl{}technology-assisted review\textquotedbl{} OR \textquotedbl{}citation screening\textquotedbl{})\par}
\item[q1b, support query] {\footnotesize\ttfamily\raggedright (\textquotedbl{}when to stop\textquotedbl{} OR \textquotedbl{}stop screening\textquotedbl{} OR \textquotedbl{}early stopping\textquotedbl{}) AND (screening) AND (\textquotedbl{}systematic review\textquotedbl{} OR \textquotedbl{}technology-assisted review\textquotedbl{})\par}
\item[q1c, support query] {\footnotesize\ttfamily\raggedright (\textquotedbl{}active learning\textquotedbl{} OR \textquotedbl{}screening prioritisation\textquotedbl{} OR \textquotedbl{}screening prioritization\textquotedbl{} OR \textquotedbl{}prioritised screening\textquotedbl{}) AND (\textquotedbl{}systematic review\textquotedbl{} OR \textquotedbl{}citation screening\textquotedbl{}) AND (recall OR \textquotedbl{}work saved\textquotedbl{})\par}
\item[q1d, support query] {\footnotesize\ttfamily\raggedright (\textquotedbl{}capture-recapture\textquotedbl{} OR \textquotedbl{}work saved over sampling\textquotedbl{} OR \textquotedbl{}recall estimation\textquotedbl{}) AND (screening OR stopping) AND (\textquotedbl{}systematic review\textquotedbl{} OR \textquotedbl{}technology-assisted review\textquotedbl{})\par}
\item[q1e, support query] {\footnotesize\ttfamily\raggedright (\textquotedbl{}technology assisted review\textquotedbl{} OR \textquotedbl{}technology-assisted review\textquotedbl{} OR \textquotedbl{}machine learning\textquotedbl{} OR \textquotedbl{}active learning\textquotedbl{} OR \textquotedbl{}prioritised screening\textquotedbl{} OR \textquotedbl{}prioritized screening\textquotedbl{} OR \textquotedbl{}screening prioritisation\textquotedbl{} OR \textquotedbl{}screening prioritization\textquotedbl{}) AND (stopping OR stop OR \textquotedbl{}target recall\textquotedbl{} OR \textquotedbl{}recall target\textquotedbl{} OR \textquotedbl{}high recall\textquotedbl{} OR \textquotedbl{}total recall\textquotedbl{}) AND (\textquotedbl{}systematic review\textquotedbl{} OR \textquotedbl{}systematic reviews\textquotedbl{} OR \textquotedbl{}citation screening\textquotedbl{} OR \textquotedbl{}abstract screening\textquotedbl{})\par}
\item[q1f, support query] {\footnotesize\ttfamily\raggedright (\textquotedbl{}machine learning\textquotedbl{} OR \textquotedbl{}active learning\textquotedbl{} OR \textquotedbl{}semi-automated\textquotedbl{} OR \textquotedbl{}semi-automation\textquotedbl{} OR \textquotedbl{}automated screening\textquotedbl{} OR \textquotedbl{}screening prioritization\textquotedbl{} OR \textquotedbl{}screening prioritisation\textquotedbl{} OR \textquotedbl{}relevance prediction\textquotedbl{} OR \textquotedbl{}relevance predictions\textquotedbl{}) AND (\textquotedbl{}title and abstract screening\textquotedbl{} OR \textquotedbl{}abstract screening\textquotedbl{} OR \textquotedbl{}citation screening\textquotedbl{} OR \textquotedbl{}title screening\textquotedbl{}) AND (\textquotedbl{}systematic review\textquotedbl{} OR \textquotedbl{}systematic reviews\textquotedbl{} OR \textquotedbl{}rapid review\textquotedbl{}) AND (sensitivity OR recall OR missed OR \textquotedbl{}relevant studies\textquotedbl{} OR workload)\par}
\item[q1g, support query] {\footnotesize\ttfamily\raggedright (prioritization OR prioritisation OR prioritized OR prioritised OR filtering OR ranking) AND (\textquotedbl{}study selection\textquotedbl{} OR \textquotedbl{}title and abstract screening\textquotedbl{} OR \textquotedbl{}abstract screening\textquotedbl{} OR \textquotedbl{}citation screening\textquotedbl{}) AND (\textquotedbl{}systematic review\textquotedbl{} OR \textquotedbl{}systematic reviews\textquotedbl{}) AND (\textquotedbl{}work savings\textquotedbl{} OR \textquotedbl{}work saved\textquotedbl{} OR recall OR sensitivity OR missed OR stopping)\par}
\item[q1x, counter query] {\footnotesize\ttfamily\raggedright (\textquotedbl{}stopping rule\textquotedbl{} OR \textquotedbl{}stopping rules\textquotedbl{} OR \textquotedbl{}stopping criteria\textquotedbl{} OR \textquotedbl{}stopping method\textquotedbl{}) AND (unreliable OR fail OR missed OR premature OR undershoot OR unstable) AND (screening) AND (review)\par}
\item[q1y, counter query] {\footnotesize\ttfamily\raggedright (\textquotedbl{}stopping rule\textquotedbl{} OR \textquotedbl{}stopping rules\textquotedbl{} OR \textquotedbl{}stopping criteria\textquotedbl{} OR \textquotedbl{}stopping criterion\textquotedbl{} OR \textquotedbl{}stopping method\textquotedbl{} OR \textquotedbl{}when to stop\textquotedbl{} OR \textquotedbl{}early stopping\textquotedbl{} OR \textquotedbl{}active learning\textquotedbl{} OR \textquotedbl{}technology-assisted review\textquotedbl{} OR \textquotedbl{}screening prioritisation\textquotedbl{} OR \textquotedbl{}screening prioritization\textquotedbl{}) AND (missed OR miss OR fail OR failed OR unreliable OR overestimate OR \textquotedbl{}below target\textquotedbl{} OR \textquotedbl{}not reached\textquotedbl{} OR variability) AND (recall) AND (\textquotedbl{}systematic review\textquotedbl{} OR \textquotedbl{}systematic reviews\textquotedbl{} OR \textquotedbl{}citation screening\textquotedbl{})\par}
\end{description}


\begin{table}
  \begin{sidecaption}{%
      Hits per query and database for the first question (stopping a
      ranked screen early), searched 2026-11-23 with Scholar, as
      \texttt{scholar sessions show} reports them; the queries are
      listed in \cref{sec:stopping-method}.
      OpenAlex (OA; not open access) in title and abstract, Semantic
      Scholar (S2) through its bulk search, Scopus as
      \texttt{TITLE-ABS-KEY(\dots)}, Web of Science (WoS) as
      \texttt{TS=(\dots)}, IEEE Xplore (IEEE), arXiv; no limit, so the
      counts are all the records each database returned.
      \enquote{113 of about 142}: the database reported more records
      than it returned (Semantic Scholar's totals are approximate).
      DBLP was not searched (\cref{sec:search-databases}).
    }[tab:stopping-matrix]
    \footnotesize
    \begin{tabular}{@{}l*{6}{>{\raggedleft\arraybackslash}p{1.45cm}}@{}}
\toprule
Query & OA & S2 & Scopus & WoS & IEEE & arXiv \\
\midrule
q1a & 71 & 38 & 35 & 25 & 2 & 6 \\
q1b & 35 & 33 & 20 & 15 & 1 & 3 \\
q1c & 58 & 45 & 47 & 26 & 6 & 4 \\
q1d & 36 & 34 & 26 & 20 & 3 & 1 \\
q1e & 193 & 105 & 116 & 73 & 37 & 11 \\
q1f & 148 & 113 of about 142 & 123 & 86 & 8 & 7 \\
q1g & 224 & 127 & 91 & 117 & 5 & 3 \\
q1x & 29 & 10 & 7 & 9 & 0 & 0 \\
q1y & 14 & 8 & 4 & 6 & 0 & 0 \\
\bottomrule
\end{tabular}

  \end{sidecaption}
\end{table}

Five studies were named before the search: two evaluations of stopping
rules (one in education, one in psychology), the statistical stopping
criterion, an analysis of the work-saved measure, and the
capture-recapture stopping rule; all five came back.
The check by meaning returned 29 works that no session held; the two
repair queries brought back 17 of them (\cref{sec:search-check}).
Of the 931 distinct records, 454 were screened in the first tier and 53
of them bore on the question (23 supporting the claim, 30 qualifying or
contradicting it); of the 477 left unscreened, a random sample of 60
held 3 bearing records (5\,\%; 95\,\% interval 1.4 to 13.7\,\%).
All three are evaluations of screening tools, published in journals of
research-synthesis methods, read only in their abstracts and not cited;
they agree with the sources below.
The records judged bearing are listed in \cref{sec:stopping-hits}.

\section{Results}
\label{sec:stopping-results}

\paragraph{Most stopping rules miss their target.}
The largest comparison found tested 15 stopping methods on 81 datasets
of real reviews and found that
\textcquote[abstract]{Repke2026}{almost all existing stopping methods
either fail to reliably stop without missing relevant records or fail to
utilize the full potential work-savings. Only one method reliably meets
the set recall target, but stops conservatively.}
An independent test of the stop criteria offered with one widely used
open tool, in 35,000 simulated screenings of five health datasets,
reports that \textcquote[abstract]{Kempny2026}{\textins*{n}one of the
tested stop criteria could consistently identify all relevant studies
across all datasets}; the share of a dataset that had to be screened to
find all relevant studies ranged from 2.9\,\% to 76.9\,\% on average.
A heuristic that stops after 50 irrelevant records in a row missed its
target of 95\,\% recall in 39\,\% of the cases in the simulations of
\textcite[p.~11]{Callaghan2020}.
A live mental-health review, reported in a preprint read only in its
abstract, concludes that \textcquote[abstract]{Slade2026}{stopping rules
risk missing relevant studies}.
For a ranked web source, a study read only in its abstract found that
\textcquote[abstract]{Briscoe2023}{the lowest ranking relevant studies
were ranked 227th and 215th for the two systematic reviews}, so the first
100 results of Google Search would have missed them.

\paragraph{A statistical criterion holds its target, at a cost.}
The one kind of rule that held a stated recall target in these
comparisons tests, on a random sample of the records not yet seen,
whether the target has been reached:
\textcquote[abstract]{Callaghan2020}{The stopping criteria are shown on
test datasets to achieve a reliable level of recall, while still
providing work reductions of on average 17\,\%. Other methods proposed
previously are shown to provide inconsistent recall and work reductions
across datasets.}
In ten health reviews, stopping a prioritised screen at the point where
95\,\% of the relevant records had in fact been found cut the screening
by a median of 47.1\,\%, and \textcquote[abstract]{Hamel2020}{\textins*{a}mong
the 100 simulations (10 simulations per review), none of these 5\,\% of
records were a final included study in the systematic review}; but that
stopping point was known only afterwards, not estimated during the
screen.
Using a tool's relevance predictions in place of one of two reviewers
missed a median of 1\,\% (0 to 14\,\%) of relevant records in eleven
systematic reviews, and in six rapid reviews
\textcquote[abstract]{Gates2020}{\textins*{t}he median (range)
proportion missed records for both approaches was 6 (0--22)\,\%}.

\paragraph{How well a rule does depends on the dataset.}
In 27 reviews in education and educational psychology, heuristic rules
retrieved 95\,\% of the relevant abstracts, but
\textcquote[abstract]{Campos2024}{the performance of the heuristic
stopping criteria depended on the learning algorithm used and the length
and proportion of relevant papers in an abstract collection}.
An evaluation in psychology, read only in its abstract, found
\textcquote[abstract]{Koenig2024}{performance differences between
stopping rules regarding all performance measures and variations in the
performance of stopping rules across different prevalence ratios}.
The authors of a conservative heuristic that combines several rules leave
the decision to the reviewer:
\textcquote[abstract]{Boetje2024}{the decision to stop the active
learning process depends on careful consideration of the trade-off
between the costs of additional record labeling against the potential
errors of the current model for the specific dataset and context}.
A commentary by the authors of the statistical criterion, with
methodologists of evidence synthesis, states the general point:
\textcquote[p.~2]{Callaghan2024}{\textins*{w}ithout prior knowledge of
the number of relevant records, any approach to stopping early risks
missing relevant studies that would have been retrieved had all records
been screened}.

\section{What a sample of the rest can say}
\label{sec:stopping-sample}

The statistical criterion rests on
\textcquote[p.~5]{Callaghan2020}{drawing a random sample from the
remaining unseen documents}.
The procedure of \cref{ch:procedure} uses the same kind of evidence for a
smaller purpose.
It does not stop when a recall target is reached; it screens a tier
chosen in advance, draws a seeded random sample of what the tier left
out, and reports the share of bearing records in the sample with its
interval.
That share estimates how much of the unscreened rest bears on the
question; it does not say which records those are, and a sample in which
none bears does not show that none does.
No source found tests such a sample against a full screen; the reason to
draw it is that it is the only statement about the unscreened records
that the sources above leave room for.

\section{Discussion and limitations}
\label{sec:stopping-discussion}

The evidence is mostly simulation on collections of reviews already
screened in full, most of them in health research; one evaluation is in
education, one in psychology, and one live review is reported only in a
preprint's abstract.
The one result in which no included study was lost at the stopping point
used a point known only in retrospect.
None of the sources measures whether the studies missed at a stopping
point would have changed a review's conclusions.
The psychology evaluation and the study of Google Search could be read
only in their abstracts (the publisher refused the full texts), and the
capture-recapture stopping rule, one of the named studies, was not
obtained and is not cited.
The search itself had one screener, an agent (\cref{sec:search-limits}).

\section{Conclusion}
\label{sec:stopping-conclusion}

No: in the comparisons found, most stopping rules either missed relevant
records or saved little work, and only a statistical criterion built on
a random sample of the unseen records held its recall target, saving
about a sixth of the screening on test collections; how well any rule
did depended on the dataset.
For the procedure this means that a tier is not a stopping rule.
It decides what is read first; the records it leaves out are counted and
sampled, and the method says so (\cref{ch:procedure}).

\section{Future work}
\label{sec:stopping-future}

This search left undone the full texts of the psychology evaluation and
the study of Google Search, and the capture-recapture rule, which a
reader with access to the publishers can add.
The field lacks a test of stopping rules outside health and education,
in computing in particular, and a measurement of whether the studies a
rule misses change what a review concludes; either would say how much a
tier can be trusted to stand in for a full screen.

\section{The records judged bearing}
\label{sec:stopping-hits}

% Author's note: table by generate_audit_table on a copy of screenlit-q1,
% bearing rows only, cited = cited in doc/bibliography.bib.
% --- scholar LaTeX fragment --------------------------------------------
% Generated by: scholar sessions export -f table
% Session: screenlit-q1
% \input this file from the body of a document whose
% preamble loads:
%   \usepackage{longtable}
%   \usepackage{booktabs}
%   \usepackage{array}
% ----------------------------------------------------------------------

\begingroup\footnotesize
\begin{longtable}{@{}%
  >{\raggedright\arraybackslash}p{0.44\textwidth}c%
  >{\raggedright\arraybackslash}p{0.13\textwidth}%
  >{\raggedright\arraybackslash}p{0.26\textwidth}@{}}
\caption{Records bearing on the claim, Question 1, stopping a ranked screen early (11 cited, 22 supporting and 20 qualifying or contradicting it, out of 897 unique hits; the 165 adjacent hits, the 282 false hits and 397 unclassified are not listed). The last column gives the reason for each inclusion; for rows decided by an agent the agent's confidence is given. A row is cited when the text cites the work, whoever decided on its records. Databases searched: OpenAlex (OA; not open access), IEEE Xplore (IEEE), Web of Science (WoS), Scopus, Semantic Scholar (S2), arXiv. Searched 2026-11-23.}\label{tab:screenlit-q1}\\
\toprule
Title & Year & Database & Reason \\
\midrule
\endfirsthead
\multicolumn{4}{@{}l}{\emph{\tablename~\thetable{} (cont.)}}\\
\toprule
Title & Year & Database & Reason \\
\midrule
\endhead
\midrule \multicolumn{4}{r@{}}{\emph{continued on next page}}\\
\endfoot
\bottomrule
\endlastfoot
An evaluation of DistillerSR's machine learning-based prioritization tool for title/abstract screening -- impact on reviewer-relevant outcomes & 2020 & OA, WoS, Scopus, S2 & \emph{cited}: supports the claim (0.70) \\
An evaluation of the performance of stopping rules in AI-aided screening for psychological meta-analytical research & 2024 & Scopus, WoS & \emph{cited}: qualifies or contradicts the claim (0.70) \\
Computer-assisted screening in systematic evidence synthesis requires robust and well-evaluated stopping criteria & 2024 & OA, Scopus & \emph{cited}: qualifies or contradicts the claim (0.60) \\
Don't Stop Me Now, `Cause I'm Having a Good Time Screening: Evaluation of Stopping Methods for Safe Use of Priority Screening in Systematic Reviews & 2026 & OA, S2 & \emph{cited}: qualifies or contradicts the claim (0.75) \\
Feasibility and desirability of screening search results from Google Search exhaustively for systematic reviews: A crosscase analysis & 2023 & OA, WoS, Scopus, S2 & \emph{cited}: qualifies or contradicts the claim (0.60) \\
Machine learning-assisted versus conventional title and abstract screening for a mental health systematic review (Preprint) & 2026 & OA & \emph{cited}: qualifies or contradicts the claim (0.70) \\
Screening Smarter, Not Harder: A Comparative Analysis of Machine Learning Screening Algorithms and Heuristic Stopping Criteria for Systematic Reviews in Educational Research & 2024 & OA, Scopus, S2, WoS & \emph{cited}: qualifies or contradicts the claim (0.70) \\
Statistical stopping criteria for automated screening in systematic reviews & 2020 & OA, WoS, Scopus, S2 & \emph{cited}: qualifies or contradicts the claim (0.80) \\
The SAFE procedure: a practical stopping heuristic for active learning-based screening in systematic reviews and meta-analyses & 2024 & WoS, Scopus, S2, OA & \emph{cited}: qualifies or contradicts the claim (0.60) \\
The semi-automation of title and abstract screening: a retrospective exploration of ways to leverage Abstrackr's relevance predictions in systematic and rapid reviews & 2020 & OA, S2, WoS, Scopus & \emph{cited}: qualifies or contradicts the claim (0.60) \\
When to stop reviewing: validation of stop criteria in ASReview & 2026 & OA, Scopus, WoS, S2 & \emph{cited}: qualifies or contradicts the claim (0.70) \\
A text-mining tool generated title-abstract screening workload savings: performance evaluation versus single-human screening & 2022 & OA, S2, WoS, Scopus & supports the claim (0.70) \\
AI-supported screening for systematic literature reviews informing EULAR recommendations and points to consider & 2026 & OA, S2, Scopus & supports the claim (0.75) \\
Can artificial intelligence separate the wheat from the chaff in systematic reviews of health economic articles? & 2023 & OA, WoS, Scopus, S2 & supports the claim (0.55) \\
Confidence-Based Stopping Methods for Systematic Reviews & 2026 & arXiv, WoS, OA, S2, Scopus & supports the claim (0.65) \\
Decision-Theoretic Stopping Rules for Document Screening & 2026 & OA, arXiv, WoS, Scopus, S2 & supports the claim (0.55) \\
Enhancing recall in automated record screening: A resampling algorithm & 2024 & OA, WoS, Scopus, S2 & supports the claim (0.50) \\
Improving Stopping Methods for Technology Assisted Review & 2026 & OA & supports the claim (0.55) \\
Increasing the Efficiency of High-Recall Information Retrieval & 2019 & OA, S2 & supports the claim (0.55) \\
Increasing the efficiency of study selection for systematic reviews using prioritization tools and a single-screening approach & 2023 & OA, S2, WoS, Scopus & supports the claim (0.60) \\
Machine learning enables automated screening for systematic reviews and meta-analysis in urology & 2024 & OA, Scopus, S2 & supports the claim (0.50) \\
Machine Learning Reduced Workload for COVID-19 Literature & 2025 & OA & supports the claim (0.50) \\
Natural language processing was effective in assisting rapid title and abstract screening when updating systematic reviews & 2021 & WoS, Scopus & supports the claim (0.60) \\
Research Screener: a machine learning tool to semi-automate abstract screening for systematic reviews & 2021 & OA, S2, WoS, Scopus & supports the claim (0.60) \\
RLStop: A Reinforcement Learning Stopping Method for TAR & 2024 & arXiv, WoS & supports the claim (0.60) \\
Sensitivity and specificity of alternative screening methods for systematic reviews using text mining tools & 2023 & WoS, Scopus & supports the claim (0.55) \\
Stopping Methods for Technology-assisted Reviews Based on Point Processes & 2023 & OA, arXiv, WoS, Scopus, S2 & supports the claim (0.65) \\
SWIFT-Active Screener: Accelerated document screening through active learning and integrated recall estimation & 2020 & OA, WoS, Scopus, S2 & supports the claim (0.65) \\
Text mining to support abstract screening for knowledge syntheses: a semi-automated workflow & 2021 & OA, S2, WoS, Scopus & supports the claim (0.55) \\
The capture--mark--recapture technique can be used as a stopping rule when searching in systematic reviews & 2008 & OA, WoS, Scopus, S2 & supports the claim (0.45) \\
Usefulness of machine learning softwares to screen titles of systematic reviews: a methodological study & 2023 & OA, S2, WoS, Scopus & supports the claim (0.50) \\
Using Chao's Estimator as a Stopping Criterion for Technology-Assisted Review & 2024 & arXiv, OA, WoS, Scopus, S2 & supports the claim (0.60) \\
Validity of using a semi-automated screening tool in a systematic review assessing non-specific effects of respiratory vaccines & 2025 & OA, S2, WoS, Scopus & supports the claim (0.50) \\
A Study on a Stopping Strategy for Systematic Reviews Based on a Distributed Effort Approach & 2020 & OA, WoS, Scopus & qualifies or contradicts the claim (0.45) \\
Application of Machine Learning Classifiers in Rapid Reviews for Health Research: A Case Example Using EPPIReviewer & 2026 & OA & qualifies or contradicts the claim (0.45) \\
Assessing the Accuracy of Machine-Assisted Abstract Screening With DistillerAI: A User Study & 2019 & OA & qualifies or contradicts the claim (0.65) \\
Assessing the accuracy of machine-assisted abstract screening with DistillerAI: a user study & 2019 & OA, S2, WoS, Scopus & qualifies or contradicts the claim (0.60) \\
Can machine learning help accelerate article screening for systematic reviews? Yes, when article separability in embedding space is high & 2025 & Scopus & qualifies or contradicts the claim (0.65) \\
Decoding semi-automated title-abstract screening: findings from a convenience sample of reviews & 2020 & OA, S2, WoS, Scopus & qualifies or contradicts the claim (0.55) \\
Faster title and abstract screening? Evaluating Abstrackr, a semi-automated online screening program for systematic reviewers & 2015 & OA, S2, WoS, Scopus & qualifies or contradicts the claim (0.55) \\
Machine learning for screening prioritization in systematic reviews: Comparative performance of Abstrackr and EPPI-Reviewer & 2020 & Scopus & qualifies or contradicts the claim (0.50) \\
Measuring the impact of screening automation on meta-analyses of diagnostic test accuracy & 2019 & OA, WoS, Scopus, S2 & qualifies or contradicts the claim (0.55) \\
Optimising screening efficiency in evidence synthesis on health Technology: A simulation study using ASReview & 2026 & WoS, Scopus, OA, S2 & qualifies or contradicts the claim (0.60) \\
Performance and Usability of Machine Learning for Screening in Systematic Reviews: A Comparative Evaluation of Three Tools & 2019 & OA, S2, WoS, Scopus & qualifies or contradicts the claim (0.55) \\
Performance of Two AI Approaches in ASReview Compared With Manual Screening for Dementia Care Literature Screening: Comparative Analysis & 2026 & OA, WoS, Scopus, S2 & qualifies or contradicts the claim (0.70) \\
Prioritising references for systematic reviews with RobotAnalyst: A user study & 2018 & OA, Scopus, S2, WoS & qualifies or contradicts the claim (0.55) \\
SnowBallSLR: saturation-aware and reproducible citation searching for systematic reviews & 2026 & OA, Scopus & qualifies or contradicts the claim (0.45) \\
Strategizing AI utilization for psychological literature screening: A comparative analysis of machine learning algorithms and key factors to consider & 2026 & Scopus, OA, S2, WoS & qualifies or contradicts the claim (0.50) \\
Technology-assisted title and abstract screening for systematic reviews: a retrospective evaluation of the Abstrackr machine learning tool & 2018 & OA, S2, WoS, Scopus & qualifies or contradicts the claim (0.50) \\
The efficiency of machine learning-assisted platform for article screening in systematic reviews in orthopaedics & 2022 & OA, WoS, Scopus, S2 & qualifies or contradicts the claim (0.50) \\
To include or not to include? A prescription from the pharmacy on how to use active learning-assisted screening in systematic reviews & 2026 & OA, WoS, Scopus, S2 & qualifies or contradicts the claim (0.45) \\
Unsupervised title and abstract screening for systematic review: a retrospective case-study using topic modelling methodology & 2023 & OA, S2, WoS, Scopus & qualifies or contradicts the claim (0.50) \\
Write your abstracts carefully---The impact of abstract reporting quality on findability by semi-automated title-abstract screening tools & 2025 & OA, S2, WoS, Scopus & qualifies or contradicts the claim (0.55) \\
\end{longtable}
\endgroup


\chapter{Shortened searches and screens}
\label{ch:abbreviated}

% Author's note (verification and reproduction): session screenlit-q2
% (and screenlit-q2-top); queries q2a-q2e, q2x-q2z; check by meaning
% checks/q2*.json; samples tier2-sample-q2.tsv (seed 20261125) and
% sample60-q2.txt (seed 20261123); see the block at the top of
% search.tex.  Sources: doc/bibliography.bib, keys Ewald2020 Marshall2019
% Hartling2016 Halladay2015 vanEnst2014 Affengruber2020 Pham2016
% Smela2024 Corso2026 Wong2017 Pena2024 Barnish2019 Waffenschmidt2019
% Gartlehner2020 Stoll2019.

A search with a limit per database, a search of fewer databases, a
screen of only part of the records, and a screen by one person instead
of two all save work.
The procedure of \cref{ch:procedure} takes none of the first two
shortcuts and takes the third only with a sample of the rest; for the
fourth it asks for a second screen where the result matters.
This chapter asks what the shortcuts cost when the same reviews are done
both ways: \emph{does an abbreviated search or screen lose few included
studies and leave pooled estimates and conclusions unchanged?}
The claim is tested as a general one; a shortcut that is harmless for
some questions and not for others does not meet it.

\section{Method}
\label{sec:abbreviated-method}

The search is the one described in \cref{ch:lit-search}.
The second question's queries are listed below, and
\cref{tab:abbreviated-matrix} gives their hits per database.
The support queries are q2a to q2d, the counter queries q2x to q2z; q2e
was added after the check by meaning, with the words of methodological
shortcuts, meta-epidemiological comparisons and the value of a second
reviewer.
\begin{description}[style=nextline,font=\normalfont\bfseries]
\item[q2a, support query] {\footnotesize\ttfamily\raggedright (\textquotedbl{}rapid review\textquotedbl{} OR \textquotedbl{}rapid reviews\textquotedbl{} OR abbreviated OR streamlined) AND (\textquotedbl{}systematic review\textquotedbl{} OR \textquotedbl{}systematic reviews\textquotedbl{}) AND (compared OR comparison OR concordance OR agreement)\par}
\item[q2b, support query] {\footnotesize\ttfamily\raggedright (\textquotedbl{}single screening\textquotedbl{} OR \textquotedbl{}single-reviewer screening\textquotedbl{} OR \textquotedbl{}one reviewer\textquotedbl{} OR \textquotedbl{}single reviewer\textquotedbl{}) AND (screening) AND (double OR dual OR \textquotedbl{}two reviewers\textquotedbl{}) AND (\textquotedbl{}systematic review\textquotedbl{})\par}
\item[q2c, support query] {\footnotesize\ttfamily\raggedright (\textquotedbl{}fewer databases\textquotedbl{} OR \textquotedbl{}reduced number of databases\textquotedbl{} OR \textquotedbl{}MEDLINE only\textquotedbl{} OR \textquotedbl{}restricting databases\textquotedbl{} OR \textquotedbl{}limited search\textquotedbl{} OR \textquotedbl{}single database\textquotedbl{}) AND (search) AND (\textquotedbl{}meta-analysis\textquotedbl{} OR \textquotedbl{}systematic review\textquotedbl{}) AND (conclusions OR \textquotedbl{}effect estimates\textquotedbl{} OR \textquotedbl{}pooled estimates\textquotedbl{} OR \textquotedbl{}included studies\textquotedbl{})\par}
\item[q2d, support query] {\footnotesize\ttfamily\raggedright (\textquotedbl{}abbreviated search\textquotedbl{} OR \textquotedbl{}abbreviated searches\textquotedbl{} OR \textquotedbl{}restricted search\textquotedbl{} OR \textquotedbl{}search restrictions\textquotedbl{} OR \textquotedbl{}limited search\textquotedbl{} OR \textquotedbl{}fewer databases\textquotedbl{} OR \textquotedbl{}number of databases\textquotedbl{} OR \textquotedbl{}single database\textquotedbl{}) AND (\textquotedbl{}systematic review\textquotedbl{} OR \textquotedbl{}systematic reviews\textquotedbl{} OR \textquotedbl{}meta-analysis\textquotedbl{} OR \textquotedbl{}meta-analyses\textquotedbl{}) AND (conclusions OR \textquotedbl{}effect estimate\textquotedbl{} OR \textquotedbl{}effect estimates\textquotedbl{} OR \textquotedbl{}pooled estimate\textquotedbl{} OR \textquotedbl{}results changed\textquotedbl{} OR \textquotedbl{}change the results\textquotedbl{})\par}
\item[q2e, support query] {\footnotesize\ttfamily\raggedright (\textquotedbl{}methodological shortcuts\textquotedbl{} OR shortcuts OR \textquotedbl{}rapid review\textquotedbl{} OR \textquotedbl{}rapid reviews\textquotedbl{} OR expedite OR expedited OR \textquotedbl{}second reviewer\textquotedbl{}) AND (\textquotedbl{}systematic review\textquotedbl{} OR \textquotedbl{}systematic reviews\textquotedbl{}) AND (\textquotedbl{}meta-epidemiological\textquotedbl{} OR \textquotedbl{}different results\textquotedbl{} OR \textquotedbl{}case studies\textquotedbl{} OR \textquotedbl{}value of\textquotedbl{} OR \textquotedbl{}same conclusions\textquotedbl{} OR \textquotedbl{}results and conclusions\textquotedbl{})\par}
\item[q2x, counter query] {\footnotesize\ttfamily\raggedright (\textquotedbl{}additional databases\textquotedbl{} OR \textquotedbl{}other databases\textquotedbl{} OR \textquotedbl{}single database\textquotedbl{} OR \textquotedbl{}one database\textquotedbl{}) AND (changed OR altered OR missed OR \textquotedbl{}would have missed\textquotedbl{} OR insufficient) AND (\textquotedbl{}systematic review\textquotedbl{} OR \textquotedbl{}meta-analysis\textquotedbl{})\par}
\item[q2y, counter query] {\footnotesize\ttfamily\raggedright (\textquotedbl{}single screening\textquotedbl{} OR \textquotedbl{}single reviewer\textquotedbl{} OR \textquotedbl{}one reviewer\textquotedbl{}) AND (missed OR missing OR \textquotedbl{}error rate\textquotedbl{} OR \textquotedbl{}false negatives\textquotedbl{}) AND (screening)\par}
\item[q2z, counter query] {\footnotesize\ttfamily\raggedright (\textquotedbl{}rapid review\textquotedbl{} OR \textquotedbl{}rapid reviews\textquotedbl{} OR \textquotedbl{}abbreviated search\textquotedbl{} OR abbreviated OR \textquotedbl{}single screening\textquotedbl{} OR \textquotedbl{}single reviewer\textquotedbl{} OR \textquotedbl{}fewer databases\textquotedbl{}) AND (\textquotedbl{}systematic review\textquotedbl{} OR \textquotedbl{}systematic reviews\textquotedbl{} OR \textquotedbl{}meta-analysis\textquotedbl{}) AND (\textquotedbl{}different conclusions\textquotedbl{} OR \textquotedbl{}changed conclusions\textquotedbl{} OR \textquotedbl{}conclusions changed\textquotedbl{} OR discordant OR \textquotedbl{}missed studies\textquotedbl{} OR \textquotedbl{}missing studies\textquotedbl{} OR \textquotedbl{}would have missed\textquotedbl{})\par}
\end{description}


\begin{table}
  \begin{sidecaption}{%
      Hits per query and database for the second question (shortened
      searches and screens), searched 2026-11-23 with Scholar, as
      \texttt{scholar sessions show} reports them; the databases and
      their syntax as in \cref{tab:stopping-matrix}.
      \enquote{of about}: Semantic Scholar's approximate total;
      \enquote{274 of 275}: the database reported one record more than
      it returned.
    }[tab:abbreviated-matrix]
    \footnotesize
    \begin{tabular}{@{}l*{6}{>{\raggedleft\arraybackslash}p{1.45cm}}@{}}
\toprule
Query & OA & S2 & Scopus & WoS & IEEE & arXiv \\
\midrule
q2a & 907 & 538 of about 541 & 479 & 509 of 510 & 11 & 6 \\
q2b & 332 & 222 of about 226 & 211 & 331 of 332 & 0 & 0 \\
q2c & 322 & 191 of about 192 & 113 & 66 & 1 & 0 \\
q2d & 469 & 336 of about 338 & 274 of 275 & 180 & 3 & 0 \\
q2e & 220 & 321 of about 322 & 165 & 79 & 0 & 1 \\
q2x & 619 & 510 of about 514 & 275 & 343 & 2 & 0 \\
q2y & 197 & 132 & 80 & 73 & 4 & 4 \\
q2z & 42 & 22 & 21 & 16 & 0 & 0 \\
\bottomrule
\end{tabular}

  \end{sidecaption}
\end{table}

Five studies were named before the search (a review of single against
double screening, two reviews of rapid-review methods, and two studies of
what sources beyond PubMed or MEDLINE add); all came back.
The check by meaning returned 36 works no session held; the repair query
returned three of the comparisons cited below and left 33 outside, most
of them general texts on systematic reviews (\cref{sec:search-check}).
Of the 4,123 distinct records, 679 were screened in the first tier, and
38 of them bore on the question (12 supporting the claim, 26 qualifying
or contradicting it).
Of the 3,384 records left unscreened, none among a random sample of 60
bore on the question (95\,\% interval 0 to 6.0\,\%); nor did any of a
second sample of 60 drawn from the records that three or more databases
had returned.
The records judged bearing are listed in \cref{sec:abbreviated-hits}.

\section{Results}
\label{sec:abbreviated-results}

\paragraph{Fewer databases: usually the same answer, not always.}
A meta-epidemiological study redid the searches of 47 Cochrane reviews in
abbreviated forms and found that the abbreviated searches led to
identical effect estimates in 34 to 79\,\% of the meta-analyses and to
opposite effects, or none that could be estimated, in 6 to 13\,\%:
\textcquote[abstract]{Ewald2020}{Abbreviated literature searches often
led to identical or very similar effect estimates as comprehensive
searches with slightly increased confidence intervals. Relevant
deviations may occur.}
In 129 Cochrane reviews from three review groups, restricting the
meta-analyses to the studies two databases held changed little:
\textcquote[abstract]{Hartling2016}{Results of meta-analyses based on the
majority of studies did not differ in most cases. There were very few
cases of changes in statistical significance.}
Studies read only in their abstracts agree: in 50 Cochrane reviews,
\textcquote[abstract]{Halladay2015}{\textins*{m}eta-analysis using only
PubMed-indexed vs.\ all available studies led to a different conclusion
in a single case}; in ten reviews of diagnostic accuracy,
\textcquote[abstract]{vanEnst2014}{\textins*{r}estricting to studies
indexed on MEDLINE did not influence the summary estimates of the
meta-analyses in our sample}; and in three rapid reviews of health
technologies, searching fewer databases would have missed ten references,
yet \textcquote[abstract]{Wong2017}{nine out of the ten references were
found to have no or minimal impact on the overall findings of the
reviews}.
The largest simulation, read only in its abstract, applied rapid-review
shortcuts to 2,512 Cochrane reviews; statistical significance changed in
6.5 to 38.6\,\% of the meta-analyses, with no evidence of bias, and its
authors conclude that \textcquote[abstract]{Marshall2019}{\textins*{s}earching
PubMed only might be considered where a \textasciitilde10\,\% risk of
the primary outcome OR changing by \textgreater20\,\% could be
tolerated}.

\paragraph{Broad questions lose more.}
Abbreviated searches combined with screening by one reviewer gave the
same conclusions as the Cochrane reviews for two oncological topics, but
for a public-health topic the rapid reviews found 38 to 48\,\% of the
included studies, and the authors conclude:
\textcquote[abstract]{Affengruber2020}{For complex public health topics,
sensitivity seems to be insufficient.}
The other comparisons outside focused clinical questions were read only
in their abstracts.
In three reviews in agri-food public health, the shortcuts usually
omitted at least one relevant study, and
\textcquote[abstract]{Pham2016}{\textins*{t}he omission of studies
affected 39 of 143 possible meta-analyses, of which 14 were no longer
possible because of insufficient studies}.
Of three expedited reviews checked against their Cochrane versions, one
missed 3 of 19 studies, and \textcquote[abstract]{Smela2024}{their lack
impacted significantly the final results, or the possibility to run
meta-analyses, in four cases}; a second missed 3 of 27 without changing
the results, and the third matched completely.
Restricting three reviews of low back pain to the three databases the
Cochrane rapid-review guidance names left out studies held by Chinese and
physiotherapy databases: \textcquote[abstract]{Corso2026}{Relative recall
ranged from 65.4\,\% to 100.0\,\%, missing up to 34.6\,\% of relevant
studies.}
In an epidemiological review, a conference abstract reports that
\textcquote[abstract]{Barnish2019}{114 studies were eligible for
inclusion, of which SR identified 101 (89\,\%, due to absence of forward
citation chasing) and RR 64 (56\,\%, McNemar test p\textless0.001)}.
The one comparison found in computing redid a software-engineering
review as a rapid review: its search returned 78\,\% of the papers the
systematic review included, and \textcquote[abstract]{Pena2024}{\textins*{b}oth
secondary studies identified the same number of pieces of evidence (30),
but the pieces of evidence are not the same}.

\section{One screener or two}
\label{sec:abbreviated-single}

A methodological review of single against double screening found that
\textcquote[abstract]{Waffenschmidt2019}{\textins*{t}he median proportion
of missed studies was 5\,\% (range 0 to 58\,\%)}; in 3 of 7 single
screenings whose effect on the findings could be assessed, all by one
less experienced reviewer, the findings of the meta-analyses would have
changed substantially, while the missed studies of the other four had no
or a negligible effect.
In a randomised trial read only in its abstract,
\textcquote[abstract]{Gartlehner2020}{single-reviewer abstract screening
missed 13\,\% of relevant studies (sensitivity: 86.6\,\%; 95\,\%
confidence interval [CI], 80.6\,\%--91.2\,\%)}, against 3\,\% for two
reviewers.
In one large review, also read only in its abstract,
\textcquote[abstract]{Stoll2019}{an additional 6.6\,\% to 9.1\,\% of
eligible studies were identified during the title/abstract stage by using
two reviewers}.
What the rapid-review guidance makes of this, and how it asks for
agreement to be checked, is in \cref{ch:reporting}.

\section{Discussion and limitations}
\label{sec:abbreviated-discussion}

Almost all of this evidence is from health research, and most of it from
reviews of clinical interventions, where the question is focused and a
few databases index the field well.
The comparisons that found larger losses (public health, agri-food,
epidemiology, low back pain, software engineering) were, apart from one,
read only in their abstracts, and so were the largest simulation, the
two studies of sources beyond PubMed and MEDLINE, the trial and the
single review on second reviewers: their full texts were behind bot
checks, refused by the publisher, or not tried.
No source tests a limit on the number of results per database as such,
which is the shortcut that \cref{sec:measured-limit} measured; the
nearest is the study of a ranked web source in \cref{ch:stopping}.

\section{Conclusion}
\label{sec:abbreviated-conclusion}

Not in general: for focused clinical questions, searching fewer
databases usually left pooled estimates identical or close, but changed
them in a minority of meta-analyses; for broad questions and outside
clinical medicine the shortcuts lost a large share of the included
studies; and one screener missed more relevant studies than two.
For the procedure this means searching every live database without a
limit, saying what was not screened, and, where a result matters,
screening a second time.

\section{Future work}
\label{sec:abbreviated-future}

This search left eleven of the fifteen sources read only in their
abstracts; their full texts would give the conditions of each loss.
The field lacks a comparison of the full and the shortened search in
computing and education beyond the one software-engineering case, and a
test of a limit per database on what a review concludes; if such a test
found that a limit rarely changes conclusions in these fields, the
procedure could spend less on searching, and if not, the case for no
limit would rest on more than one review.

\section{The records judged bearing}
\label{sec:abbreviated-hits}

% Author's note: table by generate_audit_table on a copy of screenlit-q2,
% bearing rows only, cited = cited in doc/bibliography.bib.
% --- scholar LaTeX fragment --------------------------------------------
% Generated by: scholar sessions export -f table
% Session: screenlit-q2
% \input this file from the body of a document whose
% preamble loads:
%   \usepackage{longtable}
%   \usepackage{booktabs}
%   \usepackage{array}
% ----------------------------------------------------------------------

\begingroup\footnotesize
\begin{longtable}{@{}%
  >{\raggedright\arraybackslash}p{0.44\textwidth}c%
  >{\raggedright\arraybackslash}p{0.13\textwidth}%
  >{\raggedright\arraybackslash}p{0.26\textwidth}@{}}
\caption{Records bearing on the claim, Question 2, abbreviated searches and screens (16 cited, 7 supporting and 17 qualifying or contradicting it, out of 3,922 unique hits; the 82 adjacent hits, the 670 false hits and 3,130 unclassified are not listed). The last column gives the reason for each inclusion; for rows decided by an agent the agent's confidence is given. A row is cited when the text cites the work, whoever decided on its records. Databases searched: OpenAlex (OA; not open access), IEEE Xplore (IEEE), Web of Science (WoS), Scopus, Semantic Scholar (S2), arXiv. Searched 2026-11-23.}\label{tab:screenlit-q2}\\
\toprule
Title & Year & Database & Reason \\
\midrule
\endfirsthead
\multicolumn{4}{@{}l}{\emph{\tablename~\thetable{} (cont.)}}\\
\toprule
Title & Year & Database & Reason \\
\midrule
\endhead
\midrule \multicolumn{4}{r@{}}{\emph{continued on next page}}\\
\endfoot
\bottomrule
\endlastfoot
Abbreviated and comprehensive literature searches led to identical or very similar effect estimates: a meta-epidemiological study & 2020 & WoS, Scopus, OA, S2 & \emph{cited}: supports the claim (0.80) \\
An evaluation of DistillerSR's machine learning-based prioritization tool for title/abstract screening -- impact on reviewer-relevant outcomes & 2020 & OA, WoS, Scopus, S2 & \emph{cited}, pending, unclassified \\
Assessing the coverage of PubMed, Embase, OpenAlex, and Semantic Scholar for automated single-database searches in living guideline evidence surveillance: a case study of the international polycystic ovary syndrome guidelines 2023 & 2025 & Scopus & \emph{cited}, pending, unclassified \\
Assessing the validity of a rapid review against a systematic literature review. A comparison of systematic literature reviews done by Cochrane with rapid reviews and the impact on meta-analyses results & 2024 & OA, WoS, Scopus, S2 & \emph{cited}: qualifies or contradicts the claim (0.70) \\
Combining abbreviated literature searches with single-reviewer screening: three case studies of rapid reviews & 2020 & OA, WoS, Scopus, S2 & \emph{cited}: qualifies or contradicts the claim (0.75) \\
Comparing the Efficacy of Rapid Review With a Systematic Review in the Software Engineering Field & 2024 & OA, WoS, Scopus, S2 & \emph{cited}: qualifies or contradicts the claim (0.70) \\
Does a rapid review version of a large epidemiological systematic review fail to identify many eligible studies, and what implications does this have for the results of the review? & 2019 & OA, S2 & \emph{cited}: qualifies or contradicts the claim (0.60) \\
Evaluating the proportion of retrieved citations and relevant studies when database selection is restricted to MEDLINE, CENTRAL, and Embase: part 4 of a meta-research study of rapid review methodological trade-offs & 2026 & OA, WoS, Scopus, S2 & \emph{cited}: qualifies or contradicts the claim (0.70) \\
Implications of applying methodological shortcuts to expedite systematic reviews: three case studies using systematic reviews from agrifood public health & 2016 & OA, S2, WoS, Scopus & \emph{cited}: qualifies or contradicts the claim (0.75) \\
Meta-epidemiologic analysis indicates that MEDLINE searches are sufficient for diagnostic test accuracy systematic reviews & 2014 & OA, WoS, Scopus, S2 & \emph{cited}: supports the claim (0.70) \\
Rapid reviews may produce different results to systematic reviews: a meta-epidemiological study & 2018 & OA, S2, WoS, Scopus & \emph{cited}: qualifies or contradicts the claim (0.80) \\
Single-reviewer abstract screening missed 13 percent of relevant studies: a crowd-based, randomized controlled trial & 2020 & OA, WoS, Scopus, S2 & \emph{cited}: qualifies or contradicts the claim (0.75) \\
The contribution of databases to the results of systematic reviews: a cross-sectional study & 2016 & OA, WoS, Scopus, S2 & \emph{cited}: supports the claim (0.70) \\
The Impact Of Searching Fewer Databases In Health Technology Assessment Rapid Reviews & 2017 & OA, S2 & \emph{cited}: supports the claim (0.65) \\
The value of a second reviewer for study selection in systematic reviews & 2019 & OA, S2, WoS, Scopus & \emph{cited}: qualifies or contradicts the claim (0.70) \\
Using data sources beyond PubMed has a modest impact on the results of systematic reviews of therapeutic interventions & 2015 & Scopus & \emph{cited}: supports the claim (0.75) \\
Empirical Testing of Alternative Search Methods to Retrieve Utility Values for Health Economic Modelling & 2024 & OA, WoS, Scopus, S2 & supports the claim (0.50) \\
Excluding non-English publications from evidence-syntheses did not change conclusions: a meta-epidemiological study & 2019 & OA, S2, WoS, Scopus & supports the claim (0.60) \\
Increasing the efficiency of study selection for systematic reviews using prioritization tools and a single-screening approach & 2023 & OA, Scopus, WoS, S2 & supports the claim (0.55) \\
Limited search strategies were effective in finding relevant nonrandomized studies & 2006 & Scopus, S2, OA, WoS & supports the claim (0.50) \\
Limits in the search date for rapid reviews of diagnostic test accuracy studies & 2022 & OA, WoS, Scopus, S2 & supports the claim (0.60) \\
Simulated changes in systematic review screening sensitivity and specificity associated with varying co-screening proportions & 2026 & OA, Scopus, S2 & supports the claim (0.55) \\
Successful incorporation of single reviewer assessments during systematic review screening: development and validation of sensitivity and work-saved of an algorithm that considers exclusion criteria and count & 2021 & OA, WoS, Scopus, S2 & supports the claim (0.55) \\
A retrospective comparison of systematic reviews with same-topic rapid reviews & 2017 & OA, WoS, Scopus, S2 & qualifies or contradicts the claim (0.60) \\
Characteristics and recovery methods of studies falsely excluded during literature screening---a systematic review & 2022 & OA, WoS, Scopus, S2 & qualifies or contradicts the claim (0.45) \\
Cochrane systematic reviews in acupuncture: Methodological diversity in database searching & 2005 & WoS, OA, S2, Scopus & qualifies or contradicts the claim (0.45) \\
Comparison of a full systematic review versus rapid review approaches to assess a newborn screening test for tyrosinemia type 1 & 2017 & OA, WoS, Scopus, S2 & qualifies or contradicts the claim (0.60) \\
Database selection in systematic reviews: an insight through clinical neurology & 2017 & S2, OA, WoS, Scopus & qualifies or contradicts the claim (0.50) \\
Error rates of human reviewers during abstract screening in systematic reviews & 2020 & OA, WoS, Scopus, S2 & qualifies or contradicts the claim (0.55) \\
Evaluating the accuracy and impact of methodological trade-offs used in rapid reviews on results and certainty of the evidence: meta-research studies in systematic reviews of chronic low back pain & 2026 & OA & qualifies or contradicts the claim (0.65) \\
Evaluating the accuracy of full text and serial screening by one vs two independent reviewers: part 2 of a meta-research study of rapid review methodological trade-offs & 2026 & OA, WoS, Scopus, S2 & qualifies or contradicts the claim (0.60) \\
Evaluating the accuracy of screening titles and abstracts for progressing to full-text screening by one vs two independent reviewers: part 1 of a meta-research study of rapid review methodological trade-offs & 2026 & OA, WoS, Scopus, S2 & qualifies or contradicts the claim (0.65) \\
Evaluating the impact of methodological trade-offs on meta-analyses and certainty of the evidence according to GRADE: part 5 of a meta-research study of rapid reviews & 2026 & OA, WoS, Scopus, S2 & qualifies or contradicts the claim (0.60) \\
Examining screening test accuracy in Tyrosinemia type 1: A comparison of rapid and systematic review Methods & 2016 & OA, S2 & qualifies or contradicts the claim (0.50) \\
Identification of randomized controlled trials in systematic reviews: accuracy and reliability of screening records & 2002 & OA, WoS, Scopus, S2 & qualifies or contradicts the claim (0.60) \\
Lessons for search strategies from a systematic review, in The Cochrane Library, of nutritional supplementation trials in patients after hip fracture & 2001 & WoS, OA, Scopus, S2 & qualifies or contradicts the claim (0.55) \\
Literature search in systematic reviews: How much is good enough? & 2023 & OA, WoS, Scopus & qualifies or contradicts the claim (0.60) \\
Literature searching for social science systematic reviews: consideration of a range of search techniques & 2010 & OA, WoS, Scopus, S2 & qualifies or contradicts the claim (0.50) \\
Network meta-analyses could be improved by searching more sources and by involving a librarian & 2014 & OA, S2, WoS, Scopus & qualifies or contradicts the claim (0.45) \\
Sources of evidence for systematic reviews of interventions in diabetes & 2005 & WoS, OA, Scopus & qualifies or contradicts the claim (0.55) \\
\end{longtable}
\endgroup


\chapter{Databases and the head of a ranked list}
\label{ch:databases}

% Author's note (verification and reproduction): session screenlit-q3
% (287 MB; loaded once, for the matrix and the table) and screenlit-q3-top;
% queries q3a-q3h, q3x, q3y; check by meaning checks/q3*.json; samples
% tier2-sample-q3.tsv (seed 20261126) and sample60-q3.txt (seed
% 20261123); see the block at the top of search.tex.  Sources:
% doc/bibliography.bib, keys Bramer2017 Bramer2016 Aagaard2016 Gehanno2013
% Gusenbauer2020 Haddaway2015 Sampson2006 Heck2024 Wanyama2021 Rajit2025
% Ali2022 Hartling2016.

The procedure of \cref{ch:procedure} searches every database Scholar
offers that answers, and screens the first 100 places of each answer ordered by
relevance.
Both choices assume something about databases: that one or two would not
do, and that the head of a ranked list is the right place to start but
not a place to stop.
This chapter asks the question in the form that would make both choices
unnecessary: \emph{do a few bibliographic databases, or the first
results of a ranked source, find nearly all the relevant records a
review needs?}
\enquote{Nearly all} is taken as a recall that a review could report
without a caveat.

\section{Method}
\label{sec:databases-method}

The search is the one described in \cref{ch:lit-search}.
The third question's queries are listed below, and
\cref{tab:databases-matrix} gives their hits per database.
The support queries are q3a to q3f, the counter queries q3x and q3y; q3g
was added after the check by meaning, with the words of studies that
compare what different databases find, and q3h because two of the named
studies had not come back on the question's concepts.
\begin{description}[style=nextline,font=\normalfont\bfseries]
\item[q3a, support query] {\footnotesize\ttfamily\raggedright (\textquotedbl{}Google Scholar\textquotedbl{} OR Scopus OR \textquotedbl{}Web of Science\textquotedbl{}) AND (coverage OR overlap OR \textquotedbl{}unique records\textquotedbl{} OR \textquotedbl{}unique references\textquotedbl{}) AND (\textquotedbl{}systematic review\textquotedbl{} OR \textquotedbl{}bibliographic databases\textquotedbl{})\par}
\item[q3b, support query] {\footnotesize\ttfamily\raggedright (OpenAlex OR \textquotedbl{}Semantic Scholar\textquotedbl{}) AND (coverage OR overlap OR comparison OR completeness) AND (database OR bibliographic)\par}
\item[q3c, support query] {\footnotesize\ttfamily\raggedright (\textquotedbl{}Google Scholar\textquotedbl{}) AND (\textquotedbl{}systematic review\textquotedbl{} OR \textquotedbl{}evidence review\textquotedbl{}) AND (\textquotedbl{}first 100\textquotedbl{} OR \textquotedbl{}first 200\textquotedbl{} OR \textquotedbl{}first 50\textquotedbl{} OR \textquotedbl{}top 100\textquotedbl{} OR \textquotedbl{}first pages\textquotedbl{} OR \textquotedbl{}first 1000\textquotedbl{})\par}
\item[q3d, support query] {\footnotesize\ttfamily\raggedright (\textquotedbl{}optimal combination\textquotedbl{} OR \textquotedbl{}database combination\textquotedbl{} OR \textquotedbl{}number of databases\textquotedbl{} OR \textquotedbl{}which databases\textquotedbl{} OR \textquotedbl{}additional database\textquotedbl{}) AND (\textquotedbl{}systematic review\textquotedbl{} OR \textquotedbl{}literature search\textquotedbl{}) AND (recall OR unique OR coverage)\par}
\item[q3e, support query] {\footnotesize\ttfamily\raggedright (\textquotedbl{}software engineering\textquotedbl{} OR \textquotedbl{}computer science\textquotedbl{} OR computing OR education OR \textquotedbl{}educational research\textquotedbl{}) AND (\textquotedbl{}digital libraries\textquotedbl{} OR \textquotedbl{}digital library\textquotedbl{} OR \textquotedbl{}bibliographic databases\textquotedbl{} OR \textquotedbl{}search engines\textquotedbl{} OR Scopus OR \textquotedbl{}Web of Science\textquotedbl{} OR \textquotedbl{}Google Scholar\textquotedbl{}) AND (\textquotedbl{}systematic literature review\textquotedbl{} OR \textquotedbl{}systematic review\textquotedbl{} OR \textquotedbl{}literature search\textquotedbl{}) AND (coverage OR overlap OR recall OR \textquotedbl{}search strategy\textquotedbl{} OR completeness)\par}
\item[q3f, support query] {\footnotesize\ttfamily\raggedright (\textquotedbl{}relevance ranking\textquotedbl{} OR \textquotedbl{}relevance-ranked\textquotedbl{} OR \textquotedbl{}first 100\textquotedbl{} OR \textquotedbl{}first 200\textquotedbl{} OR \textquotedbl{}first 300\textquotedbl{} OR \textquotedbl{}top results\textquotedbl{} OR \textquotedbl{}first pages\textquotedbl{}) AND (\textquotedbl{}Google Scholar\textquotedbl{} OR \textquotedbl{}search engine\textquotedbl{} OR \textquotedbl{}web search\textquotedbl{} OR \textquotedbl{}academic search engine\textquotedbl{}) AND (\textquotedbl{}systematic review\textquotedbl{} OR \textquotedbl{}literature search\textquotedbl{} OR \textquotedbl{}grey literature\textquotedbl{})\par}
\item[q3g, support query] {\footnotesize\ttfamily\raggedright (Scopus OR \textquotedbl{}Web of Science\textquotedbl{} OR \textquotedbl{}Google Scholar\textquotedbl{} OR EBSCO OR \textquotedbl{}bibliographic databases\textquotedbl{}) AND (\textquotedbl{}systematic literature review\textquotedbl{} OR \textquotedbl{}systematic review\textquotedbl{} OR \textquotedbl{}systematic reviews\textquotedbl{} OR \textquotedbl{}meta-analyses\textquotedbl{}) AND (\textquotedbl{}common to\textquotedbl{} OR shared OR variation OR \textquotedbl{}choice of\textquotedbl{} OR \textquotedbl{}stand-alone\textquotedbl{} OR \textquotedbl{}single database\textquotedbl{} OR sufficient)\par}
\item[q3h, support query] {\footnotesize\ttfamily\raggedright (\textquotedbl{}multiple databases\textquotedbl{} OR \textquotedbl{}number of databases\textquotedbl{} OR \textquotedbl{}additional databases\textquotedbl{} OR \textquotedbl{}selective database searching\textquotedbl{} OR \textquotedbl{}diminishing returns\textquotedbl{} OR \textquotedbl{}added value\textquotedbl{}) AND (\textquotedbl{}systematic review\textquotedbl{} OR \textquotedbl{}systematic reviews\textquotedbl{} OR \textquotedbl{}meta-analysis\textquotedbl{}) AND (searching OR search OR retrieval OR recall)\par}
\item[q3x, counter query] {\footnotesize\ttfamily\raggedright (\textquotedbl{}Google Scholar\textquotedbl{}) AND (inadequate OR \textquotedbl{}not suitable\textquotedbl{} OR limitations OR \textquotedbl{}not recommended\textquotedbl{} OR \textquotedbl{}cannot be used\textquotedbl{}) AND (\textquotedbl{}systematic review\textquotedbl{} OR searching)\par}
\item[q3y, counter query] {\footnotesize\ttfamily\raggedright (\textquotedbl{}Google Scholar\textquotedbl{} OR Scopus OR \textquotedbl{}Web of Science\textquotedbl{} OR OpenAlex OR \textquotedbl{}Semantic Scholar\textquotedbl{} OR \textquotedbl{}bibliographic databases\textquotedbl{}) AND (coverage OR overlap OR \textquotedbl{}unique records\textquotedbl{} OR recall) AND (missed OR missing OR incomplete OR insufficient OR \textquotedbl{}not sufficient\textquotedbl{} OR \textquotedbl{}not indexed\textquotedbl{}) AND (\textquotedbl{}systematic review\textquotedbl{} OR \textquotedbl{}systematic reviews\textquotedbl{} OR \textquotedbl{}literature search\textquotedbl{})\par}
\end{description}


\begin{table}
  \begin{sidecaption}{%
      Hits per query and database for the third question (databases and
      ranked lists), searched 2026-11-23 with Scholar, as \texttt{scholar
      sessions show} reports them; the databases and their syntax as in
      \cref{tab:stopping-matrix}.
      \enquote{of}: the database reported more records than it returned
      (\enquote{of about}: Semantic Scholar's approximate total).
      \textsuperscript{a}~Incomplete; Scholar recorded \enquote{rate
      limited or HTTP/parse error}, and the search's log showed rate
      limits (HTTP 429).
      \textsuperscript{b}~Incomplete; Scopus stopped with HTTP 400 after
      4,991 records.
      The OpenAlex answer to q3g holds 9,978 of the 26,809 records
      OpenAlex reported, and Scholar did not mark it incomplete.
    }[tab:databases-matrix]
    \footnotesize
    \begin{tabular}{@{}l*{6}{>{\raggedleft\arraybackslash}p{1.45cm}}@{}}
\toprule
Query & OA & S2 & Scopus & WoS & IEEE & arXiv \\
\midrule
q3a & 6,237 & 999 of about 5,207\textsuperscript{a} & 4,266 of 4,272 & 3,574 of 3,576 & 50 & 11 \\
q3b & 581 & 227 of about 228 & 118 & 103 & 6 & 0 \\
q3c & 204 & 104 & 101 & 57 & 0 & 0 \\
q3d & 43 & 59 & 43 & 19 & 0 & 0 \\
q3e & 3,973 & 2,971 of about 3,008 & 1,721 of 1,723 & 1,084 of 1,086 & 149 & 51 \\
q3f & 477 of 479 & 116 & 97 & 59 & 0 & 0 \\
q3g & 9,978 of 26,809 & 7,996 of about 47,668\textsuperscript{a} & 4,991 of 18,141\textsuperscript{b} & 14,448 of 14,458 & 172 & 34 \\
q3h & 6,825 of 6,838 & 999 of about 5,845\textsuperscript{a} & 3,867 of 3,870 & 4,519 of 4,521 & 21 & 2 \\
q3x & 8,296 & 9,956 of about 10,680\textsuperscript{a} & 4,924 of 4,927 & 6,178 of 6,181 & 107 & 71 \\
q3y & 1,001 & 1,279 of about 1,294 & 391 & 374 of 375 & 9 & 4 \\
\bottomrule
\end{tabular}

  \end{sidecaption}
\end{table}

Five studies were named before the search: two on database combinations,
one on the contribution of databases to the results of reviews, one on
Google Scholar in evidence reviews, and one older study of searching
several databases.
The study of the contribution of databases came back only through the
second question's queries, and the older study through no query; it is
not cited.
The check by meaning returned 29 works no session held; the repair query
returned one of them, and 26 comparisons of bibliographic coverage stayed
outside (\cref{sec:search-check}).
Of the 61,348 distinct records, 1,003 were screened in the first tier and
17 more that five or more databases had returned; 58 bore on the question
(22 supporting the claim, 36 qualifying or contradicting it).
Of the 60,268 records left unscreened, none among a random sample of 60
bore on the question (95\,\% interval 0 to 6.0\,\%), and none of a second
sample of 60 from the records that three or more databases had returned.
The records judged bearing are listed in \cref{sec:databases-hits}.

\section{Results}
\label{sec:databases-results}

\paragraph{No single database's search, and few added ones, suffice in health.}
In 58 health reviews, \textcquote[abstract]{Bramer2017}{\textins*{s}ixteen
percent of the included references (291 articles) were only found in a
single database; Embase produced the most unique references (n =
132)}.
Embase, MEDLINE, Web of Science and the first 200 results of Google
Scholar together reached a recall of 98.3\,\%, and the authors estimate
that 60\,\% of published reviews miss more than 5\,\% of the relevant
references because of the databases they searched.
In 23 reviews of musculoskeletal disorders, more databases added little:
\textcquote[abstract]{Aagaard2016}{\textins*{c}umulative median recall for
combined searching in MEDLINE, EMBASE and CENTRAL was 88.9\,\% and
increased to 90.9\,\% when adding 10 additional databases}.

\paragraph{Holding a record is not returning it.}
All 738 studies included in 29 medical reviews were in Google Scholar
when looked up by title: \textcquote[abstract]{Gehanno2013}{All the 738
original studies included in the gold standard database were retrieved
in GS (100\,\%).}
In 120 reviews, Google Scholar held 97.2\,\% of the included references,
but a search returned far fewer near the top:
\textcquote[abstract]{Bramer2016}{\textins*{o}nly 46.4\,\% of the included
references were among the downloadable first 1000 references in GS}.
A test of 28 search systems, read only in its abstract, concludes that
\textcquote[abstract]{Gusenbauer2020}{our findings demonstrate why Google
Scholar is inappropriate as principal search system}.
For OpenAlex and Semantic Scholar the one study found measured holding,
not returning: it looked the articles a health guideline included up by
identifier, then by title.
Read only in its abstract, it reports that \textcquote[abstract]{Rajit2025}{OpenAlex
had the best coverage (98.6\,\%), followed by Semantic Scholar (98.3\,\%),
Embase (96.8\,\%), and PubMed (93.0\,\%)}.

\paragraph{The head of a ranked list.}
A search engine ranked the Boolean results of nine Cochrane reviews by
relevance; the recall of included studies within the first 500 places
was 0.70, and \textcquote[abstract]{Sampson2006}{\textins*{t}he relevance
ranking provided by the search engine was better than expected by
chance}, though it could not tell apart records that reviewers had
already pre-screened.
In environmental science, Google Scholar missed important literature in
five of six case studies, and the authors recommend reading further than
the first page: \textcquote[abstract]{Haddaway2015}{If used in systematic
reviews for grey literature, we recommend that searches of article titles
focus on the first 200 to 300 results.}
The study of Google Search in \cref{ch:stopping} found relevant studies at
places 215 and 227.

\paragraph{Outside health.}
\begingroup\sloppy
In three reviews in education, \textcquote[abstract]{Heck2024}{discipline-\hspace{0pt}specific
databases outperform international multidisciplinary sources, and a
combination of discipline-\hspace{0pt}specific international and national sources is
most efficient in finding a high proportion of relevant literature}.
In a social-science review, read only in its abstract,
\textcquote[abstract]{Wanyama2021}{only 130 articles (5.7\,\% of 2267
retrieved) were found common to all three databases}.
For reference and citation searching in software engineering, one case
(\textcquote[abstract]{Ali2022}{\textins*{i}n this particular case}) read
only in its abstract suggests that \textcquote[abstract]{Ali2022}{Scopus
seems to be more effective and useful for snowball sampling than GS for systematic
secondary studies attempting to identify peer-reviewed literature}.
How much the database choice changes a review's results, rather than its
records, is in \cref{ch:abbreviated}; the cross-sectional study of 129
Cochrane reviews there \parencite{Hartling2016} is also one of this
question's named studies.
\par\endgroup

\section{Discussion and limitations}
\label{sec:databases-discussion}

The recall figures are from health research and from databases Scholar
does not search (Embase, MEDLINE, CENTRAL, Google Scholar); they show that
no single source's search sufficed there, not which combination suffices in
computing or education.
In education the discipline's own databases did best, and Scholar
searches none of them; a claim about education should say so.
No study found tests the keyword or ranked retrieval of OpenAlex or
Semantic Scholar; the one that measures them counts records held by
identifier, and was read only in its abstract.
Our own measurements are the only evidence here about the head of
OpenAlex, arXiv and IEEE Xplore lists, and they are one chapter's
sessions (\cref{sec:measured-chapter}).
The search's own answers were large and partly incomplete
(\cref{tab:databases-matrix}), so for this question the first tier and the
sample stand for more unscreened records than for the others.

\section{Conclusion}
\label{sec:databases-conclusion}

No: in health research no single database's search found what the
reviews included, three or four together came close, and further ones
added
little; holding a record is not returning it, since a ranked source that
held almost every included study returned fewer than half of them in its
first 1,000 places; and outside health the right databases differ.
For the procedure this means searching every database that answers,
reading the head of each ranked list first without stopping there, and
not carrying the health figures over to computing or education.

\section{Future work}
\label{sec:databases-future}

This search did not obtain the full texts of the test of 28 search
systems, of the study of OpenAlex and Semantic Scholar, of the
social-science and software-engineering comparisons, or the older study
of searching several databases.
The field lacks a measurement of the recall of OpenAlex's and Semantic
Scholar's keyword searches for the studies reviews include, and of the
databases computing reviews need; such a study would say whether the
six databases of \cref{ch:procedure} are enough or which one to add.

\section{The records judged bearing}
\label{sec:databases-hits}

The first table lists the bearing records decided in the question's
session; the second those decided in its companion session of the first
100 places, which holds seven bearing records the first table lacks (the
cited works appear in both).

% Author's note: tables by generate_audit_table on copies of screenlit-q3
% and screenlit-q3-top, bearing rows only, cited = cited in
% doc/bibliography.bib.  Five of the seven bearing records of the -top
% session are pending in screenlit-q3; two (arXiv) are not in it.
% --- scholar LaTeX fragment --------------------------------------------
% Generated by: scholar sessions export -f table
% Session: screenlit-q3
% \input this file from the body of a document whose
% preamble loads:
%   \usepackage{longtable}
%   \usepackage{booktabs}
%   \usepackage{array}
% ----------------------------------------------------------------------

\begingroup\footnotesize
\begin{longtable}{@{}%
  >{\raggedright\arraybackslash}p{0.44\textwidth}c%
  >{\raggedright\arraybackslash}p{0.13\textwidth}%
  >{\raggedright\arraybackslash}p{0.26\textwidth}@{}}
\caption{Records bearing on the claim, Question 3, databases and ranked lists (15 cited, 17 supporting and 23 qualifying or contradicting it, out of 60,049 unique hits; the 158 adjacent hits, the 814 false hits and 59,022 unclassified are not listed). The last column gives the reason for each inclusion; for rows decided by an agent the agent's confidence is given. A row is cited when the text cites the work, whoever decided on its records. Databases searched: OpenAlex (OA; not open access), IEEE Xplore (IEEE), Web of Science (WoS), Scopus (incomplete for 1 of 10 queries), Semantic Scholar (S2; incomplete for 4 of 10 queries), arXiv. Searched 2026-11-23.}\label{tab:screenlit-q3}\\
\toprule
Title & Year & Database & Reason \\
\midrule
\endfirsthead
\multicolumn{4}{@{}l}{\emph{\tablename~\thetable{} (cont.)}}\\
\toprule
Title & Year & Database & Reason \\
\midrule
\endhead
\midrule \multicolumn{4}{r@{}}{\emph{continued on next page}}\\
\endfoot
\bottomrule
\endlastfoot
A Comparison of Citation Sources for Reference and Citation-Based Search in Systematic Literature Reviews & 2022 & WoS, Scopus & \emph{cited}: qualifies or contradicts the claim (0.60) \\
Assessing the coverage of PubMed, Embase, OpenAlex, and Semantic Scholar for automated single-database searches in living guideline evidence surveillance: a case study of the international polycystic ovary syndrome guidelines 2023 & 2025 & OA, WoS, Scopus, S2 & \emph{cited}: supports the claim (0.75) \\
Can electronic search engines optimize screening of search results in systematic reviews: an empirical study & 2006 & OA, Scopus & \emph{cited}: qualifies or contradicts the claim (0.60) \\
Comparing the coverage, recall, and precision of searches for 120 systematic reviews in Embase, MEDLINE, and Google Scholar: a prospective study & 2016 & OA, WoS, Scopus, S2 & \emph{cited}: qualifies or contradicts the claim (0.70) \\
Coverage and similarity of bibliographic databases to find most relevant literature for systematic reviews in education & 2024 & WoS, OA, Scopus, S2 & \emph{cited}: qualifies or contradicts the claim (0.75) \\
Evaluating the proportion of retrieved citations and relevant studies when database selection is restricted to MEDLINE, CENTRAL, and Embase: part 4 of a meta-research study of rapid review methodological trade-offs & 2026 & Scopus, OA & \emph{cited}, pending, unclassified \\
Feasibility and desirability of screening search results from Google Search exhaustively for systematic reviews: A cross-case analysis & 2023 & Scopus & \emph{cited}, pending, unclassified \\
Is the coverage of google scholar enough to be used alone for systematic reviews & 2013 & WoS, Scopus, S2, OA & \emph{cited}: supports the claim (0.60) \\
Optimal database combinations for literature searches in systematic reviews: a prospective exploratory study & 2017 & OA, WoS, Scopus, S2 & \emph{cited}: supports the claim (0.80) \\
Optimizing literature search in systematic reviews - are MEDLINE, EMBASE and CENTRAL enough for identifying effect studies within the area of musculoskeletal disorders? & 2016 & Scopus, OA, WoS, S2 & \emph{cited}: supports the claim (0.60) \\
The contribution of databases to the results of systematic reviews: a cross-sectional study & 2016 & OA, WoS, Scopus & \emph{cited}, pending, unclassified \\
The Impact Of Searching Fewer Databases In Health Technology Assessment Rapid Reviews & 2017 & OA & \emph{cited}, pending, unclassified \\
The Role of Google Scholar in Evidence Reviews and Its Applicability to Grey Literature Searching & 2015 & OA, WoS, Scopus, S2 & \emph{cited}: qualifies or contradicts the claim (0.65) \\
Where you search determines what you find: the effects of bibliographic databases on systematic reviews & 2022 & WoS, OA & \emph{cited}: qualifies or contradicts the claim (0.75) \\
Which academic search systems are suitable for systematic reviews or meta-analyses? Evaluating retrieval qualities of Google Scholar, PubMed, and 26 other resources & 2020 & WoS & \emph{cited}: qualifies or contradicts the claim (0.80) \\
A Comparison of OpenAlex With Scopus and Web of Science for Tracking Scholarly Nursing Literature & 2025 & Scopus, OA, WoS, S2 & supports the claim (0.55) \\
A search of only four key databases would identify most randomized controlled trials of acupuncture: A metaepidemiological study & 2022 & OA, WoS, Scopus, S2 & supports the claim (0.70) \\
Analyzing the Utility of Openalex to Identify Studies for Systematic Reviews: Methods and a Case Study & 2025 & OA, S2 & supports the claim (0.55) \\
Can Google Scholar Serve as a Sufficient Source for Systematic Reviews? Reconsidering the Need for Multiple Database Searches & 2026 & OA & supports the claim (0.40) \\
Database coverage and their use in systematic reviews regarding spinal manipulative therapy: an exploratory study & 2022 & WoS, Scopus, OA, S2 & supports the claim (0.55) \\
Google Scholar as a Sole Literature Source for Meta-analyses in Otolaryngology Head and Neck Surgery & 2023 & Scopus, OA, S2, WoS & supports the claim (0.60) \\
Google Scholar Could Be Used as a Stand-Alone Resource for Systematic Reviews & 2015 & WoS, OA & supports the claim (0.50) \\
Google Scholar Out-Performs Many Subscription Databases when Keyword Searching & 2010 & OA, WoS & supports the claim (0.40) \\
Google Scholar, Microsoft Academic, Scopus, Dimensions, Web of Science, and OpenCitations' COCI: a multidisciplinary comparison of coverage via citations & 2021 & WoS, Scopus & supports the claim (0.45) \\
Guidelines for the Search Strategy to Update Systematic Literature Reviews in Software Engineering & 2020 & arXiv, WoS, Scopus, S2 & supports the claim (0.50) \\
Knowing When to Stop: Final Results vs. Work Involved in Systematic Review Database Searching & 2016 & OA & supports the claim (0.55) \\
On the performance of hybrid search strategies for systematic literature reviews in software engineering & 2020 & OA, arXiv, WoS, Scopus, S2 & supports the claim (0.55) \\
Reassessing the Role of Google Scholar in PRISMA-Informed Systematic Reviews Through a Critical Analysis of Influential Research Identifying Its Limitations and Empirical Evidence & 2026 & OA, WoS, Scopus, S2 & supports the claim (0.50) \\
Redefining searching in non-medical sciences systematic reviews: The ascendance of Google Scholar as the primary database & 2025 & Scopus, OA, WoS, S2 & supports the claim (0.50) \\
Reference coverage analysis of OpenAlex compared to Web of Science and Scopus & 2025 & OA, WoS, Scopus, S2 & supports the claim (0.50) \\
Sources Used To Find Studies For Systematic Reviews Of Economic Evaluations & 2017 & OA, S2 & supports the claim (0.50) \\
The Utility of Google Scholar When Searching Geographical Literature: Comparison With Three Commercial Bibliographic Databases & 2015 & OA, WoS, Scopus, S2 & supports the claim (0.40) \\
A comparison of the performance of seven key bibliographic databases in identifying all relevant systematic reviews of interventions for hypertension & 2016 & OA, WoS, Scopus, S2 & qualifies or contradicts the claim (0.55) \\
A systematic review case study of urgent and emergency care configuration found citation searching of Web of Science and Google Scholar of similar value & 2022 & OA, WoS, Scopus, S2 & qualifies or contradicts the claim (0.55) \\
A Systematic Review on the Sufficiency of PubMed and Google Scholar for Biosciences & 2020 & OA, S2 & qualifies or contradicts the claim (0.55) \\
An alternative screening approach for Google Search identifies an accurate and manageable number of results for a systematic review (case study) & 2021 & OA, WoS, Scopus, S2 & qualifies or contradicts the claim (0.55) \\
An analysis of the suitability of OpenAlex for bibliometric analyses & 2024 & OA, S2 & qualifies or contradicts the claim (0.50) \\
Beyond Google Scholar, Scopus, and Web of Science: An evaluation of the backward and forward citation coverage of 59 databases' citation indices & 2024 & Scopus, OA, WoS, S2 & qualifies or contradicts the claim (0.50) \\
Can we prioritise which databases to search? A case study using a systematic review of frozen shoulder management & 2012 & OA, WoS, Scopus, S2 & qualifies or contradicts the claim (0.55) \\
Citation searching: a systematic review case study of multiple risk behaviour interventions & 2014 & OA, WoS, Scopus, S2 & qualifies or contradicts the claim (0.50) \\
Coverage of the bibliographic databases in mental health research & 2009 & OA, WoS, Scopus, S2 & qualifies or contradicts the claim (0.45) \\
Evaluating the linguistic coverage of OpenAlex : An assessment of metadata accuracy and completeness & 2025 & OA, WoS, Scopus, S2 & qualifies or contradicts the claim (0.45) \\
Google Scholar as replacement for systematic literature searches: good relative recall and precision are not enough & 2013 & OA, WoS, Scopus, S2 & qualifies or contradicts the claim (0.50) \\
Google Scholar is not enough to be used alone for systematic reviews & 2013 & OA & qualifies or contradicts the claim (0.65) \\
How many sources are needed? The effects of bibliographic databases on systematic review outcomes & 2022 & WoS, OA, Scopus, S2, IEEE & qualifies or contradicts the claim (0.70) \\
Identifying Relevant Controlled Clinical Trials for Systematic Reviews Requires Searching Multiple Resources -- and, Even Then, Comprehensiveness is Questionable & 2006 & OA, S2 & qualifies or contradicts the claim (0.55) \\
Large coverage fluctuations in Google Scholar: A case study & 2021 & Scopus, WoS, OA, arXiv, S2 & qualifies or contradicts the claim (0.50) \\
Overlaps of multiple database retrieval and citation tracking in dementia care research: a methodological study & 2021 & OA, WoS, Scopus, S2 & qualifies or contradicts the claim (0.50) \\
Search where you will find most: Comparing the disciplinary coverage of 56 bibliographic databases & 2022 & OA, WoS, Scopus, S2 & qualifies or contradicts the claim (0.50) \\
Searching for qualitative health research required several databases and alternative search strategies: a study of coverage in bibliographic databases & 2019 & OA, WoS, Scopus, S2 & qualifies or contradicts the claim (0.60) \\
Searching the Veterinary Literature: A Comparison of the Coverage of Veterinary Journals by Nine Bibliographic Databases & 2012 & OA, WoS, Scopus, S2 & qualifies or contradicts the claim (0.45) \\
Understanding discrepancies in the coverage of OpenAlex : The case of China & 2025 & OA, WoS, Scopus, S2 & qualifies or contradicts the claim (0.45) \\
Using Conventional Bibliographic Databases for Social Science Research: Web of Science and Scopus are not the Only Options & 2021 & OA, Scopus, S2 & qualifies or contradicts the claim (0.40) \\
Variation in number of hits for complex searches in Google Scholar & 2016 & WoS, OA & qualifies or contradicts the claim (0.45) \\
Web Search Engines - Not Yet a Reliable Replacement for Bibliographic Databases & 2018 & OA, WoS, Scopus, S2 & qualifies or contradicts the claim (0.60) \\
\end{longtable}
\endgroup

% --- scholar LaTeX fragment --------------------------------------------
% Generated by: scholar sessions export -f table
% Session: screenlit-q3-top
% \input this file from the body of a document whose
% preamble loads:
%   \usepackage{longtable}
%   \usepackage{booktabs}
%   \usepackage{array}
% ----------------------------------------------------------------------

\begingroup\footnotesize
\begin{longtable}{@{}%
  >{\raggedright\arraybackslash}p{0.44\textwidth}c%
  >{\raggedright\arraybackslash}p{0.13\textwidth}%
  >{\raggedright\arraybackslash}p{0.26\textwidth}@{}}
\caption{Records bearing on the claim, Question 3, records decided in the companion session of the first 100 places (9 cited, 1 supporting and 6 qualifying or contradicting it, out of 979 unique hits; the 11 adjacent hits, the 85 false hits and 867 unclassified are not listed). The last column gives the reason for each inclusion; for rows decided by an agent the agent's confidence is given. A row is cited when the text cites the work, whoever decided on its records. Databases searched: OpenAlex (OA; not open access; capped at the limit for 6 of 10 queries), arXiv. Searched 2026-11-23.}\label{tab:screenlit-q3-top}\\
\toprule
Title & Year & Database & Reason \\
\midrule
\endfirsthead
\multicolumn{4}{@{}l}{\emph{\tablename~\thetable{} (cont.)}}\\
\toprule
Title & Year & Database & Reason \\
\midrule
\endhead
\midrule \multicolumn{4}{r@{}}{\emph{continued on next page}}\\
\endfoot
\bottomrule
\endlastfoot
Assessing the coverage of PubMed, Embase, OpenAlex, and Semantic Scholar for automated single-database searches in living guideline evidence surveillance: a case study of the international polycystic ovary syndrome guidelines 2023 & 2025 & OA & \emph{cited}, pending, unclassified \\
Can electronic search engines optimize screening of search results in systematic reviews: an empirical study & 2006 & OA & \emph{cited}, pending, unclassified \\
Comparing the coverage, recall, and precision of searches for 120 systematic reviews in Embase, MEDLINE, and Google Scholar: a prospective study & 2016 & OA & \emph{cited}, pending, unclassified \\
Coverage and similarity of bibliographic databases to find most relevant literature for systematic reviews in education & 2023 & OA & \emph{cited}, pending, unclassified \\
Is the coverage of google scholar enough to be used alone for systematic reviews & 2013 & OA & \emph{cited}, pending, unclassified \\
Optimal database combinations for literature searches in systematic reviews: a prospective exploratory study & 2017 & OA & \emph{cited}, pending, unclassified \\
Optimizing literature search in systematic reviews -- are MEDLINE, EMBASE and CENTRAL enough for identifying effect studies within the area of musculoskeletal disorders? & 2016 & OA & \emph{cited}, pending, unclassified \\
The Role of Google Scholar in Evidence Reviews and Its Applicability to Grey Literature Searching & 2015 & OA & \emph{cited}, pending, unclassified \\
Where you search determines what you find: the effects of bibliographic databases on systematic reviews & 2021 & OA & \emph{cited}, pending, unclassified \\
Complementary approaches to searching MEDLINE may be sufficient for updating systematic reviews & 2016 & OA & supports the claim (0.50) \\
Cochrane Systematic Reviews in Acupuncture: Methodological Diversity in Database Searching & 2005 & OA & qualifies or contradicts the claim (0.45) \\
Combination of multiple databases is necessary for a valid systematic review & 2014 & OA & qualifies or contradicts the claim (0.60) \\
Evaluating the linguistic coverage of OpenAlex : An assessment of metadata accuracy and completeness & 2025 & OA, arXiv & qualifies or contradicts the claim (0.40) \\
Literature search in systematic reviews: How much is good enough? & 2023 & OA & qualifies or contradicts the claim (0.60) \\
Searching Additional Databases Except PubMed Are Necessary for a Systematic Review & 2008 & OA & qualifies or contradicts the claim (0.55) \\
Understanding discrepancies in the coverage of OpenAlex : The case of China & 2025 & OA, arXiv & qualifies or contradicts the claim (0.40) \\
\end{longtable}
\endgroup


\chapter{Screening with a language model}
\label{ch:llm-screening}

% Author's note (verification and reproduction): session screenlit-q4
% (and screenlit-q4-top); queries q4a-q4e, q4x-q4z; check by meaning
% checks/q4.json; sample tier2-sample-q4.tsv (seed 20261127); see the
% block at the top of search.tex.  Sources: doc/bibliography.bib, keys
% Xie2026 Tran2024 Sandner2025 Krag2024 Khraisha2024 Sciurti2025
% Huotala2026 Figalova2026 Teo2023 Mateen2013.

Scholar's classifier screens records with a language model, and an agent
can screen them too (\cref{ch:procedure}).
A model's verdict is cheap, and a cheaper first pass still (titles only,
a small model, a model trusted where it is confident) would be cheaper.
This chapter asks what such a pass costs in missed studies:
\emph{does a first pass with a language model miss few relevant studies
compared with human screening?}
The claim is tested as one a procedure could lean on without a human
check: few missed studies, across reviews.

\section{Method}
\label{sec:llm-method}

The search is the one described in \cref{ch:lit-search}.
The fourth question's queries are listed below, and
\cref{tab:llm-matrix} gives their hits per database.
The support queries are q4a to q4e, the counter queries q4x to q4z; no
query was added after the check by meaning.
\begin{description}[style=nextline,font=\normalfont\bfseries]
\item[q4a, support query] {\footnotesize\ttfamily\raggedright (\textquotedbl{}large language model\textquotedbl{} OR \textquotedbl{}large language models\textquotedbl{} OR GPT OR ChatGPT) AND (\textquotedbl{}title and abstract\textquotedbl{} OR \textquotedbl{}title-and-abstract\textquotedbl{} OR \textquotedbl{}abstract screening\textquotedbl{} OR \textquotedbl{}citation screening\textquotedbl{}) AND (\textquotedbl{}systematic review\textquotedbl{} OR \textquotedbl{}literature review\textquotedbl{})\par}
\item[q4b, support query] {\footnotesize\ttfamily\raggedright (\textquotedbl{}title only\textquotedbl{} OR \textquotedbl{}titles only\textquotedbl{} OR \textquotedbl{}title-only\textquotedbl{} OR \textquotedbl{}title screening\textquotedbl{}) AND (screening) AND (abstract) AND (\textquotedbl{}systematic review\textquotedbl{})\par}
\item[q4c, support query] {\footnotesize\ttfamily\raggedright (\textquotedbl{}large language model\textquotedbl{} OR GPT OR LLM) AND (screening) AND (cascade OR \textquotedbl{}two-stage\textquotedbl{} OR uncertain OR \textquotedbl{}smaller model\textquotedbl{} OR cost OR \textquotedbl{}cost-effective\textquotedbl{}) AND (\textquotedbl{}systematic review\textquotedbl{})\par}
\item[q4d, support query] {\footnotesize\ttfamily\raggedright (\textquotedbl{}large language model\textquotedbl{} OR GPT OR LLM OR ChatGPT) AND (screening) AND (sensitivity OR recall OR missed OR \textquotedbl{}false negative\textquotedbl{} OR error) AND (human OR reviewers)\par}
\item[q4e, support query] {\footnotesize\ttfamily\raggedright (\textquotedbl{}language model\textquotedbl{} OR \textquotedbl{}language models\textquotedbl{} OR LLM OR LLMs OR \textquotedbl{}GPT-4\textquotedbl{} OR \textquotedbl{}zero-shot\textquotedbl{}) AND (screening OR \textquotedbl{}study selection\textquotedbl{} OR \textquotedbl{}citation screening\textquotedbl{} OR \textquotedbl{}abstract screening\textquotedbl{}) AND (\textquotedbl{}systematic review\textquotedbl{} OR \textquotedbl{}systematic reviews\textquotedbl{} OR \textquotedbl{}literature review\textquotedbl{}) AND (sensitivity OR recall OR specificity OR missed)\par}
\item[q4x, counter query] {\footnotesize\ttfamily\raggedright (LLM OR GPT OR ChatGPT) AND (screening) AND (unreliable OR inconsistent OR poor OR \textquotedbl{}not suitable\textquotedbl{} OR limitations OR \textquotedbl{}low sensitivity\textquotedbl{}) AND (\textquotedbl{}systematic review\textquotedbl{})\par}
\item[q4y, counter query] {\footnotesize\ttfamily\raggedright (\textquotedbl{}large language model\textquotedbl{} OR \textquotedbl{}large language models\textquotedbl{} OR \textquotedbl{}language model\textquotedbl{} OR LLM OR GPT OR ChatGPT) AND (\textquotedbl{}title and abstract\textquotedbl{} OR \textquotedbl{}abstract screening\textquotedbl{} OR \textquotedbl{}citation screening\textquotedbl{} OR \textquotedbl{}title screening\textquotedbl{} OR screening) AND (\textquotedbl{}systematic review\textquotedbl{} OR \textquotedbl{}systematic reviews\textquotedbl{}) AND (\textquotedbl{}low sensitivity\textquotedbl{} OR missed OR \textquotedbl{}false negatives\textquotedbl{} OR unreliable OR inconsistent OR \textquotedbl{}not replace\textquotedbl{} OR \textquotedbl{}cannot replace\textquotedbl{} OR poor)\par}
\item[q4z, counter query] {\footnotesize\ttfamily\raggedright (\textquotedbl{}title only\textquotedbl{} OR \textquotedbl{}titles only\textquotedbl{} OR \textquotedbl{}title-only\textquotedbl{} OR \textquotedbl{}title screening\textquotedbl{}) AND (screening) AND (missed OR \textquotedbl{}lower sensitivity\textquotedbl{} OR \textquotedbl{}false negatives\textquotedbl{} OR insufficient OR \textquotedbl{}relevant studies\textquotedbl{}) AND (\textquotedbl{}systematic review\textquotedbl{} OR \textquotedbl{}systematic reviews\textquotedbl{})\par}
\end{description}


\begin{table}
  \begin{sidecaption}{%
      Hits per query and database for the fourth question (screening
      with a language model), searched 2026-11-23 with Scholar, as
      \texttt{scholar sessions show} reports them; the databases and
      their syntax as in \cref{tab:stopping-matrix}.
      \enquote{of}: the database reported more records than it returned
      (\enquote{of about}: Semantic Scholar's approximate total).
    }[tab:llm-matrix]
    \footnotesize
    \begin{tabular}{@{}l*{6}{>{\raggedleft\arraybackslash}p{1.45cm}}@{}}
\toprule
Query & OA & S2 & Scopus & WoS & IEEE & arXiv \\
\midrule
q4a & 262 & 218 of about 220 & 197 & 94 & 13 & 24 \\
q4b & 247 & 202 of about 232 & 164 & 103 & 6 & 0 \\
q4c & 169 & 77 of about 78 & 88 & 43 & 9 & 16 \\
q4d & 801 & 397 of about 399 & 662 & 251 of 252 & 69 & 150 \\
q4e & 581 of 587 & 304 of about 306 & 210 & 174 & 21 & 47 \\
q4x & 196 & 136 & 116 & 85 & 14 & 19 \\
q4y & 195 & 123 & 84 & 89 & 9 & 9 \\
q4z & 51 of 52 & 44 of about 72 & 28 & 20 & 1 & 0 \\
\bottomrule
\end{tabular}

  \end{sidecaption}
\end{table}

Three studies were named before the search (an exploratory study of
screening with language models, an evaluation of GPT-4 for screening and
data extraction, and a study of automated screening for clinical
reviews); all came back.
The check by meaning returned 4 works that no session held; one was a
comparison of screening by title with screening by title and abstract,
which one query's wording should have returned, and it is cited below
although no session of this question holds it.
Of the 2,574 distinct records, 652 were screened in the first tier and 12
more that five or more databases had returned; 107 bore on the question
(63 supporting the claim, 44 qualifying or contradicting it).
Of the 1,910 records left unscreened, none among a random sample of 60
bore on the question (95\,\% interval 0 to 6.0\,\%).
The search found studies of two-stage and combined set-ups (a small model
first, a larger one on the uncertain records) but read them only in their
abstracts and cites none, so nothing is said here about such set-ups.
The search used no classifier; the screening was an agent's
(\cref{sec:search-screening}).
The records judged bearing are listed in \cref{sec:llm-hits}.

\section{Results}
\label{sec:llm-results}

\paragraph{High sensitivity on average, in health.}
A meta-analysis of 18 studies of screening medical literature with
language models, read as a preprint, reports that
\textcquote[preprint, abstract]{Xie2026}{\textins*{i}n title and abstract
screening, the pooled sensitivity was 0.92 and pooled specificity was
0.94}.
Over five reviews and 22,666 citations, for GPT-3.5 used to rule records
out, \textcquote[preprint, abstract]{Tran2024}{\textins*{p}ooled estimates
of sensitivity and specificity were 97.1\,\% (95\,\%CI 89.6\,\% to
99.2\,\%) and 37.7\,\%, (95\,\%CI 18.4\,\% to 61.9\,\%), respectively},
and limiting human screening to the records not ruled out could cut the
workload by 65\,\%.
A systematic review of such studies, read only in its abstract, found
that \textcquote[abstract]{Sandner2025}{\textins*{w}hile several
approaches achieve sensitivity above 95\,\%, none consistently reach the
99\,\% threshold required for replacing human screening}.
Small models, in a study read only in its abstract, were sensitive and
imprecise: \textcquote[abstract]{Sciurti2025}{\textins*{a}cross the
systematic reviews, LLMs achieved high sensitivity (up to 100\,\%) and
low precision (below 10\,\%) for records included by full text}, at
\$0.14 to \$1.93 per review for the model used through an interface, or
free for the two run locally.

\paragraph{Wide variation between reviews.}
On 500 abstracts of one systematic review with 12 included studies, the
best models were as accurate as human readers and still missed most of
the included studies:
\textcquote[preprint, abstract]{Krag2024}{GPT-4o, GPT-4T, and
Claude3-Opus achieved the highest accuracies (97\,\% to 98\,\%) comparable
to the human readers (96\,\% and 98\,\%), although sensitivity was low
(33\,\% to 50\,\%).}
Adjusted for chance agreement and the rarity of relevant records, GPT-4's
screening fared worse than its raw accuracy suggests:
\textcquote[abstract]{Khraisha2024}{barring studies that used highly
reliable prompts -- screening performance levelled at none to moderate
for different stages and languages}.
In software engineering, eight newer models did only slightly better than
seven older ones on a screening benchmark (mean Matthews correlation
0.365 against 0.347), differences between reviews exceeded those between
models, and the authors conclude, in a preprint, that
\textcquote[abstract]{Huotala2026}{LLMs are not yet ready to replace
humans in paper screening and the advantages new, more costly models
bring, appear to be very limited}.

\paragraph{People miss studies too, and runs differ.}
In a preregistered comparison within one complex scoping review, reported
in a preprint, two human workflows and two runs of one model reached a
similar recall of 82.3 to 82.9\,\%, none found every eligible record, and
\textcquote[abstract]{Figalova2026}{\textins*{t}wo nominally identical
GPT5.4 file-batch runs agreed on 91.7\,\% of records, yet differed on 94
records, including 29 verified eligible records retained by only one
run}.
In our own test, two identical runs of the classifier on 50 papers
differed on one kept-or-discarded verdict (\cref{sec:measured-classifier}).

\paragraph{Titles alone.}
In one nutrition review, \textcquote[abstract]{Teo2023}{57.9\,\% of
studies were correctly identified at the first-level screening stage
using the Ti/O approach versus 94.7\,\% by the Ti + Ab approach}, title
only against title and abstract.
In another review, read only in its abstract, screening titles first and
screening titles with abstracts both found every included article: \textcquote[abstract]{Mateen2013}{\textins*{t}he final
systematic review included 13 articles, all of which were identified by
both screening strategies}.

\section{Discussion and limitations}
\label{sec:llm-discussion}

Most of the evidence is from health reviews, several sources are
preprints (the meta-analysis and the study of GPT-3.5 in versions that
precede their published ones, which were not obtained), and models
change faster than studies of them appear.
Average sensitivity hides the spread between reviews, which is what a
single review faces: models as accurate as human readers missed half or
more of the included studies of one review that had few of them.
The one study in computing is a preprint.
Three sources were read only in their abstracts.
This search measured no error rate of a model of its own, since no
classifier screened it.

\section{Conclusion}
\label{sec:llm-conclusion}

Not as a general rule: on health datasets language models were sensitive
on average, but no approach reliably reached the sensitivity that
replacing human screening would need, sensitivity varied widely between
reviews, two identical runs of one model differed on eligible records,
and screening by title alone lost studies in one of two reviews.
For the procedure this means that a model's verdicts, like an agent's,
are unreviewed until a person confirms them; that a second screen and its
agreement are measured where the result matters; and that the records a
model discards are sampled like any other unscreened rest.

\section{Future work}
\label{sec:llm-future}

This search did not read in full the systematic review of sensitivity
and workload, the study of compact models, or the studies of two-stage
set-ups, and it used preprints for two sources whose published versions
exist.
The field lacks studies of model screening outside health, beyond the one
software-engineering benchmark, and of the spread between runs and
between reviews rather than averages; a study that measured both for
reviews in computing would say how large a second screen needs to be.

\section{The records judged bearing}
\label{sec:llm-hits}

% Author's note: table by generate_audit_table on a copy of screenlit-q4,
% bearing rows only, cited = cited in doc/bibliography.bib.
% --- scholar LaTeX fragment --------------------------------------------
% Generated by: scholar sessions export -f table
% Session: screenlit-q4
% \input this file from the body of a document whose
% preamble loads:
%   \usepackage{longtable}
%   \usepackage{booktabs}
%   \usepackage{array}
% ----------------------------------------------------------------------

\begingroup\footnotesize
\begin{longtable}{@{}%
  >{\raggedright\arraybackslash}p{0.44\textwidth}c%
  >{\raggedright\arraybackslash}p{0.13\textwidth}%
  >{\raggedright\arraybackslash}p{0.26\textwidth}@{}}
\caption{Records bearing on the claim, Question 4, screening with a language model (10 cited, 54 supporting and 34 qualifying or contradicting it, out of 2,478 unique hits; the 95 adjacent hits, the 507 false hits and 1,778 unclassified are not listed). The last column gives the reason for each inclusion; for rows decided by an agent the agent's confidence is given. A row is cited when the text cites the work, whoever decided on its records. Databases searched: OpenAlex (OA; not open access), IEEE Xplore (IEEE), Web of Science (WoS), Scopus, Semantic Scholar (S2), arXiv. Searched 2026-11-23.}\label{tab:screenlit-q4}\\
\toprule
Title & Year & Database & Reason \\
\midrule
\endfirsthead
\multicolumn{4}{@{}l}{\emph{\tablename~\thetable{} (cont.)}}\\
\toprule
Title & Year & Database & Reason \\
\midrule
\endhead
\midrule \multicolumn{4}{r@{}}{\emph{continued on next page}}\\
\endfoot
\bottomrule
\endlastfoot
Assessing the coverage of PubMed, Embase, OpenAlex, and Semantic Scholar for automated single-database searches in living guideline evidence surveillance: a case study of the international polycystic ovary syndrome guidelines 2023 & 2025 & Scopus & \emph{cited}, pending, unclassified \\
Can large language models replace humans in the systematic review process? Evaluating GPT-4's efficacy in screening and extracting data from peer-reviewed and grey literature in multiple languages & 2023 & arXiv, Scopus, OA, S2 & \emph{cited}: qualifies or contradicts the claim (0.60) \\
Compact large language models for title and abstract screening in systematic reviews: An assessment of feasibility, accuracy, and workload reduction & 2025 & Scopus, OA, WoS, S2 & \emph{cited}: supports the claim (0.60) \\
Evaluating human and LLM screening workflows in a conceptually complex scoping review: Recall--workload trade-offs and run-to-run consistency & 2026 & OA, arXiv, S2 & \emph{cited}: qualifies or contradicts the claim (0.60) \\
Evaluating Large Language Models for Literature Screening: A Systematic Review of Sensitivity and Workload Reduction & 2025 & OA, S2 & \emph{cited}: qualifies or contradicts the claim (0.70) \\
Has LLM Screening Performance Stalled in Software Engineering Systematic Reviews? & 2026 & arXiv, OA, S2 & \emph{cited}: qualifies or contradicts the claim (0.60) \\
Large language models for abstract screening in systematic- and scoping reviews: A diagnostic test accuracy study & 2024 & OA, S2 & \emph{cited}: qualifies or contradicts the claim (0.65) \\
Performance of Large Language Models in Automated Medical Literature Screening: A Systematic Review and Meta-analysis & 2026 & OA, WoS, Scopus & \emph{cited}: supports the claim (0.75) \\
Sensitivity and Specificity of Using GPT-3.5 Turbo Models for Title and Abstract Screening in Systematic Reviews and Meta-analyses & 2024 & Scopus, OA, WoS, S2 & \emph{cited}: supports the claim (0.70) \\
Title-plus-abstract versus title-only first-level screening approach: a case study using a systematic review of dietary patterns and sarcopenia risk to compare screening performance & 2023 & OA, WoS, Scopus, S2 & \emph{cited}: qualifies or contradicts the claim (0.70) \\
A comprehensive systematic review dataset is a rich resource for training and evaluation of AI systems for title and abstract screening & 2025 & OA, WoS, Scopus, S2 & supports the claim (0.50) \\
A Hybrid Semi-Automated Workflow for Systematic and Literature Review Processes with Large Language Model Analysis & 2024 & OA, WoS, Scopus, S2 & supports the claim (0.50) \\
A prompt-based tutorial for large language model-assisted screening in systematic reviews and meta-analyses & 2026 & WoS, Scopus, OA, S2 & supports the claim (0.50) \\
A Reproducible LLM-Ensemble Protocol for Title--Abstract Screening in OR/MS Systematic Reviews: Design and Proof of Concept & 2026 & OA, S2 & supports the claim (0.50) \\
Artificial intelligence-driven study selection in systematic reviews of randomized controlled trials, emulated trials and economic evaluation studies using large language models & 2026 & Scopus, OA, WoS, S2 & supports the claim (0.50) \\
Assessing the Ability of ChatGPT to Screen Articles for Systematic Reviews & 2023 & arXiv & supports the claim (0.50) \\
Automated full-text screening and accelerated reviews using large language models with context-aware agents: an exploratory analysis in biomarker research & 2026 & OA, WoS, Scopus, S2 & supports the claim (0.50) \\
Automated Paper Screening for Clinical Reviews Using Large Language Models: Data Analysis Study & 2024 & OA, WoS, Scopus, arXiv, S2 & supports the claim (0.50) \\
Automating Screening of Titles and Abstracts in Systematic Reviews: An Assessment of GPT-4o mini & 2026 & OA, S2 & supports the claim (0.50) \\
Collaborating with large language models in literature screening for a systematic review of college students' GenAI literacy & 2026 & OA, WoS, Scopus, S2 & supports the claim (0.50) \\
Collaborative large language models (LLMs) are all you need for screening in systematic reviews & 2026 & OA, S2 & supports the claim (0.50) \\
Comparative Assessment of ChatGPT, DeepSeek, and Human Reviewers for Full-Text Screening in Systematic Reviews on the Impact of Air Pollution in Respiratory Diseases & 2025 & OA, WoS, Scopus & supports the claim (0.50) \\
Comparing Open-Weight Large Language Models for Systematic Review Tasks: One Health Antimicrobial Resistance Case Study (Preprint) & 2026 & OA & supports the claim (0.50) \\
Cost-Efficient Large Language Model-Assisted Title and Abstract Screening for Systematic Reviews: Method Evaluation and Validation & 2026 & OA & supports the claim (0.50) \\
Development and evaluation of prompts for a large language model to screen titles and abstracts in a living systematic review & 2025 & OA, WoS, Scopus, S2 & supports the claim (0.50) \\
Development of meta-prompts for Large Language Models to screen titles and abstracts for diagnostic test accuracy reviews & 2023 & OA, S2 & supports the claim (0.50) \\
Development of Prompt Templates for Large Language Model--Driven Screening in Systematic Reviews & 2025 & Scopus, OA, WoS, S2 & supports the claim (0.50) \\
Dual-Model LLM Ensemble via Web Chat Interfaces Reaches Near-Perfect Sensitivity for Systematic-Review Screening: A Multi-Domain Validation with Equivalence to API Access & 2025 & OA, S2 & supports the claim (0.50) \\
Efficiency of LLM-Based Systematic Review Screening: Mimicking Human Expert Cognition & 2025 & OA & supports the claim (0.50) \\
Ensembles of Large Language Models for Identifying EQ-5D Studies in PubMed Based on Their Abstracts & 2026 & arXiv, OA, IEEE, Scopus, S2 & supports the claim (0.50) \\
ENTGPT: Applying Large Language Models to Systematic Review Screening With the Novel STARR Protocol & 2026 & Scopus, OA, WoS, S2 & supports the claim (0.50) \\
Evaluating a Locally Deployed 20-Billion Parameter Large Language Model for Automated Abstract Screening in Systematic Reviews & 2027 & WoS, Scopus, OA, S2 & supports the claim (0.50) \\
Evaluating large language models for title/abstract screening: a systematic review and meta-analysis \& development of new tool & 2025 & OA, WoS, Scopus, S2 & supports the claim (0.65) \\
Evaluating the Efficacy and Efficiency of GPT-5 for Automated Title and Abstract Screening in Orthopedic Surgery Systematic Reviews & 2025 & OA, WoS, Scopus, S2 & supports the claim (0.50) \\
Evaluating the Efficacy of Large Language Models for Systematic Review and Meta-Analysis Screening & 2024 & OA, S2 & supports the claim (0.50) \\
Expediting citation screening using PICo-based title-only screening for identifying studies in scoping searches and rapid reviews & 2017 & OA, WoS, Scopus, S2 & supports the claim (0.45) \\
Fine-Tuning A Large Language Model for Systematic Review Screening & 2026 & OA, arXiv, S2 & supports the claim (0.50) \\
GPT-3.5 Turbo and GPT-4 Turbo in Title and Abstract Screening for Systematic Reviews & 2025 & WoS, Scopus, OA, S2 & supports the claim (0.50) \\
High-performance automated abstract screening with large language model ensembles & 2024 & arXiv, OA, WoS, Scopus, S2 & supports the claim (0.50) \\
How Well Do ChatGPT and Claude Perform in Study Selection for Systematic Review in Obstetrics & 2025 & OA, WoS, Scopus, S2 & supports the claim (0.50) \\
Human-Comparable Sensitivity of Large Language Models in Identifying Eligible Studies Through Title and Abstract Screening: 3-Layer Strategy Using GPT-3.5 and GPT-4 for Systematic Reviews & 2024 & OA, WoS, Scopus, S2 & supports the claim (0.50) \\
Large language model-based multiagent collaboration for abstract screening toward automated systematic reviews & 2026 & WoS, Scopus, OA, S2 & supports the claim (0.50) \\
Leveraging LLMs for Title and Abstract Screening for Systematic Review: A Cost-Effective Dynamic Few-Shot Learning Approach & 2025 & arXiv, OA, WoS, Scopus, S2 & supports the claim (0.50) \\
Lightweight open-source reasoning large language model for evidence synthesis screening & 2026 & OA, WoS, Scopus, S2 & supports the claim (0.55) \\
Making Broad Evidence Synthesis Feasible: An LLM Screening Agent for Meta-Analyses Applied To Suicide Prevention & 2026 & OA, S2 & supports the claim (0.50) \\
MechaScreener: Large Language Model-Based Automated Screening for Systematic Reviews and Research & 2026 & OA, Scopus, S2 & supports the claim (0.50) \\
MechaScreener: Large Language Model-Based Automated Screening for Systematic Reviews and Research & 2026 & OA & supports the claim (0.50) \\
Methodological insights into ChatGPT's screening performance in systematic reviews & 2024 & OA, WoS, Scopus, S2 & supports the claim (0.50) \\
Operationalising LLM-assisted screening of literature to support systematic reviews & 2026 & OA, S2 & supports the claim (0.50) \\
Optimal large language models to screen citations for systematic reviews & 2025 & OA, WoS, Scopus, S2 & supports the claim (0.50) \\
Performance Comparison of Large Language Models for Efficient Literature Screening & 2025 & OA, WoS, Scopus, S2 & supports the claim (0.50) \\
Performance of a Large Language Model in Screening Citations & 2024 & Scopus, OA, WoS, S2 & supports the claim (0.50) \\
Qualitative Analysis of Discrepancy Patterns Between Large Language Models and Human Reviewers in Abstract Screening for Systematic Reviews & 2026 & OA, WoS, S2 & supports the claim (0.50) \\
REFLECTIVE-TIAB: cost-effective prompt optimization for large language model-based title and abstract screening in literature reviews & 2026 & OA, WoS, Scopus, S2 & supports the claim (0.50) \\
Synapse: Automated Generation of Systematic Review Protocols and AI-Assisted Screening with Large Language Models & 2026 & OA, S2 & supports the claim (0.50) \\
Towards Automating Articles Screening Processes Using Chain-of-Thought Large Language Models & 2026 & OA, WoS, Scopus, S2 & supports the claim (0.50) \\
Transforming literature screening: The emerging role of large language models in systematic reviews & 2025 & Scopus, S2 & supports the claim (0.50) \\
Usefulness of machine learning softwares to screen titles of systematic reviews: a methodological study & 2023 & OA, WoS, Scopus, S2 & supports the claim (0.50) \\
Using a large language model as a third reviewer to augment dual human fulltext screening in orthopaedic systematic reviews & 2026 & OA, WoS, Scopus, S2 & supports the claim (0.50) \\
Using full agreement across multiple large language models for title-and-abstract screening in systematic reviews: a proof-of-concept & 2026 & Scopus, S2 & supports the claim (0.50) \\
Using large language models to directly screen electronic databases as an alternative to traditional search strategies such as the Cochrane highly sensitive search for filtering randomized controlled trials in systematic reviews & 2025 & Scopus, WoS, OA, S2 & supports the claim (0.50) \\
Validation of Generative AI Models to Expedite Title and Abstract Screening in Systematic Reviews & 2025 & OA & supports the claim (0.50) \\
Validation of large language models (Llama 3 and ChatGPT-4o mini) for title and abstract screening in biomedical systematic reviews & 2025 & Scopus, OA, WoS, S2 & supports the claim (0.50) \\
Zero-shot Generative Large Language Models for Systematic Review Screening Automation & 2024 & arXiv, OA, WoS, Scopus, S2 & supports the claim (0.50) \\
A biomedical BERT ensemble (TITAN-SR) outperformed active-learning and large language model chatbots for systematic review title and abstract screening: development and external validation & 2026 & OA, WoS, Scopus, S2 & qualifies or contradicts the claim (0.50) \\
A Machine Learning and Large Language Model Tool for Systematic Literature Reviews of Health Economic Evidence: A Validation Study & 2026 & OA, WoS, Scopus, S2 & qualifies or contradicts the claim (0.50) \\
Automated title and abstract screening for scoping reviews using the GPT-4 Large Language Model & 2023 & OA, arXiv, S2 & qualifies or contradicts the claim (0.50) \\
Batch Size Effects on Mid2025 StateoftheArt Large Language Model Performance in Automated Title and Abstract Screening & 2026 & OA & qualifies or contradicts the claim (0.50) \\
Beyond Accuracy: LLM Variability in Evidence Screening for Software Engineering SLRs & 2026 & arXiv, OA, S2 & qualifies or contradicts the claim (0.55) \\
Cal-X: Enhancing Systematic Review Screening with LLMs and Next-Token Likelihood Calibration & 2025 & IEEE, Scopus, OA, S2 & qualifies or contradicts the claim (0.50) \\
ChatGPT (GPT-3.5) versus human reviewers for title-and-abstract screening in a scoping review: A retrospective comparative evaluation & 2026 & WoS, Scopus, OA, S2 & qualifies or contradicts the claim (0.50) \\
Class Imbalance and Batch Effects in LLM-Based Screening for Systematic Reviews & 2026 & OA, arXiv, S2 & qualifies or contradicts the claim (0.45) \\
Efficacy of Large Language Models for Screening of Systematic Reviews on Periprosthetic Joint Infection & 2026 & OA, WoS, Scopus, S2 & qualifies or contradicts the claim (0.50) \\
Evaluating and explaining LLM classification for automatic screening in Systematic Reviews & 2026 & Scopus, OA, S2 & qualifies or contradicts the claim (0.50) \\
Evaluating Large Language Models for Radiology Systematic Review Title and Abstract Screening & 2025 & OA, WoS, Scopus, S2 & qualifies or contradicts the claim (0.50) \\
Generative artificial intelligence use in evidence synthesis: A systematic review & 2025 & OA, WoS, Scopus, S2 & qualifies or contradicts the claim (0.60) \\
How Far Can a Language Model Carry a Systematic Review? Preregistered Reliability Evidence From Screening to Measurement-Quality Coding & 2026 & OA & qualifies or contradicts the claim (0.55) \\
Is Large Language Model-Assisted Citation Screening Feasible in a Scoping Review on Nonpharmacological Interventions for Delirium in Patients With Cancer? & 2025 & OA, S2 & qualifies or contradicts the claim (0.50) \\
Large Language Model-Assisted Title and Abstract Screening for A Scoping Review of Circulating miRNA Monitoring: Pilot Validation Against a Human Reviewer and Quantifying the Error of Keyword Pre-Filters & 2026 & OA & qualifies or contradicts the claim (0.50) \\
Large language models for conducting systematic reviews: on the rise, but not yet ready for use---a scoping review & 2025 & OA & qualifies or contradicts the claim (0.55) \\
Large language models for literature screening in conceptually complex and interdisciplinary reviews & 2026 & OA, WoS, Scopus, S2 & qualifies or contradicts the claim (0.50) \\
Large language models for the screening step in systematic reviews in dentistry & 2025 & OA, WoS, Scopus, S2 & qualifies or contradicts the claim (0.50) \\
Large Language Models for Title/Abstract Screening in Systematic Literature Reviews: A Case Study in Precision Livestock Farming & 2025 & OA, IEEE, Scopus, S2 & qualifies or contradicts the claim (0.50) \\
Large language models show promising performance for some systematic review tasks but call for cautious implementation: a systematic review & 2026 & OA, WoS, Scopus, S2 & qualifies or contradicts the claim (0.50) \\
Linguistic Polarity and Decision Architecture in LLM-Based Abstract Screening for Systematic Reviews & 2026 & Scopus, OA, WoS, S2 & qualifies or contradicts the claim (0.50) \\
LLM-assisted screening and extraction for systematic reviews: workload reduction and refusal-driven missingness in a land-use evidence synthesis & 2026 & OA & qualifies or contradicts the claim (0.55) \\
Performance of large language models conducting systematic review tasks in prosthodontics & 2026 & OA, WoS, Scopus, S2 & qualifies or contradicts the claim (0.50) \\
Performance of LLMs in Citation Screening: A Comparison Across Datasets with Varied Inclusion Rates & 2025 & Scopus, OA, WoS, S2 & qualifies or contradicts the claim (0.55) \\
SESR-Eval: Dataset for Evaluating LLMs in the Title-Abstract Screening of Systematic Reviews & 2025 & arXiv, OA, IEEE, WoS, Scopus, S2 & qualifies or contradicts the claim (0.50) \\
The AI Reviewer: Evaluating AI's Role in Citation Screening for Streamlined Systematic Reviews & 2025 & OA, WoS, Scopus & qualifies or contradicts the claim (0.50) \\
The Promise and Challenges of Using LLMs to Accelerate the Screening Process of Systematic Reviews & 2024 & OA, arXiv, WoS, Scopus, S2 & qualifies or contradicts the claim (0.50) \\
TiAb Review Plugin: A Browser-Based Tool for AI-Assisted Study Selection in Systematic Reviews & 2026 & arXiv, OA, S2 & qualifies or contradicts the claim (0.50) \\
Title and abstract screening for literature reviews using large language models: an exploratory study in the biomedical domain & 2024 & OA, WoS, Scopus, S2 & qualifies or contradicts the claim (0.60) \\
Understanding LLMs in Title-Abstract Screening: From Disagreements to Recommendations & 2027 & Scopus, arXiv, S2 & qualifies or contradicts the claim (0.50) \\
Using GPT4 for Title and Abstract Screening in a Literature Review of Public Policies: A Feasibility Study & 2025 & OA, S2 & qualifies or contradicts the claim (0.50) \\
Validating Large Language Models for Title-Abstract Screening in Low-Prevalence Systematic Reviews: An Environmental Science Case Study & 2026 & OA, WoS, Scopus, S2 & qualifies or contradicts the claim (0.55) \\
Validation of a two-stage automated screening pipeline for medical systematic reviews: a stratified agreement study with three independent human reviewers, using PABAK and Gwet's AC1 & 2026 & OA, S2 & qualifies or contradicts the claim (0.50) \\
Validation of an LLM-assisted title and abstract screener on two published systematic reviews & 2026 & OA & qualifies or contradicts the claim (0.60) \\
\end{longtable}
\endgroup


\chapter{What a report must say about what was not screened}
\label{ch:reporting}

% Author's note (verification and reproduction): the four guidelines were
% retrieved as known items (OpenAlex lookups, 2026-11-23) and read in full
% from their publishers' copies; the commentary was returned by the first
% question's queries.  Counter-searches: q1y and q2z (see their COUNTER
% fields).  Sources: doc/bibliography.bib, keys Page2021 Rethlefsen2021
% Garritty2024 Stevens2024 Callaghan2024.

The last step of \cref{ch:procedure} writes down what was searched, what
was screened and what was not.
This chapter asks what the reporting guidelines require of a review that
does not screen everything: \emph{what must a report say about limits on
the search, about a partial screen, and about who or what screened?}
Unlike the four chapters before it, the question is about what documents
say, so the guidelines were retrieved as known items and read in full:
the reporting statement for systematic reviews, its extension for
literature searches, and the methods and reporting guidance for rapid
reviews.
The counter queries of the first and second questions
(\cref{ch:lit-search}) returned no work that contradicts what the
guidelines say, and one commentary on what they leave out, cited below.

\section{What the guidelines ask}
\label{sec:reporting-guidelines}

\begingroup\sloppy
\begin{description}
\item[Limits on the search] Item 7 of the statement for systematic
  reviews reads: \textcquote[Table~1, item~7]{Page2021}{Present the
  full search strategies for all databases, registers and websites,
  including any filters and limits used.}
  Its extension for searches asks for the reason too, in its item 9:
  \textcquote[Table~1, item~9]{Rethlefsen2021}{Specify that no limits were
  used, or describe any limits or restrictions applied to a search (e.g.,
  date or time period, language, study design) and provide justification
  for their use.}
  It also asks for the date of each search (item 13), the number of
  records each database returned (item 15) and how duplicates were
  removed (item 16).
\item[A partial screen] Of web search engines and websites, which can
  return more results than a team can screen, the extension says:
  \textcquote[p.~180]{Rethlefsen2021}{If review teams choose to review a
  limited set of results, it should be noted in the text, along with the
  rationale.}
  The reporting guidance for rapid reviews asks authors to
  \textcquote[Table, \enquote{Distinguishing the RR from an SR}]{Stevens2024}{\textins*{e}xplicitly
  describe why the product is an RR, noting the specific steps of conduct
  or methods that characterise the distinction from an SR}.
\item[Who or what screened] Item 8 of the statement reads:
  \textcquote[Table~1, item~8]{Page2021}{Specify the methods
  used to decide whether a study met the inclusion criteria of the review,
  including how many reviewers screened each record and each report
  retrieved, whether they worked independently, and if applicable, details
  of automation tools used in the process.}
  Its flow diagram asks, where automation tools were used, how many
  records a human excluded and how many the tools excluded.
\item[Two screeners, or one] The Cochrane guidance for rapid reviews
  recommends searching at least two bibliographic databases
  (recommendation 5) and, in recommendation 9, that two reviewers screen
  at least 20\,\% of the records, check their agreement and discuss the
  discrepancies:
  \textcquote[p.~7]{Garritty2024}{Teams can proceed with single screening
  if agreement is high (eg, $\kappa$ is $\geq$0.8).}
  It adds that the 20\,\% is \textcquote[p.~7]{Garritty2024}{based on our
  experience and is therefore not evidence based}.
\end{description}
\endgroup

\section{What they leave out}
\label{sec:reporting-gaps}

No guideline found asks for the reason a screen was stopped.
A commentary on stopping criteria notes of the statement for systematic
reviews that \textcquote[p.~2]{Callaghan2024}{it only stipulates reporting
the \enquote{details of automation tools used in the process}, without
requiring a rationale for stopping decisions}.
The procedure's account of the tier, the unscreened rest and the sample
therefore goes beyond what the guidelines require; it follows the
extension's request that a limited set of results be stated with its
rationale, and applies it to every ranked answer, not only to web
search engines.

\section{What the procedure writes}
\label{sec:reporting-procedure}

Taken together, the method of a search made by \cref{ch:procedure} says:
the databases searched and those that failed, each query verbatim with
its date and its hits per database (\texttt{scholar sessions show}), that
no limit was used, how duplicates were merged; the first tier's
definition and size, the records left unscreened, and the sample of them
with its seed, size, bearing count, rate and interval (\texttt{scholar
sessions tier-sample -{}-report}); who screened, by kind (a person, a
classifier with its model, an agent) and how many of the machine's
verdicts a person reviewed; and, where a second screen was made, the
agreement (\texttt{scholar sessions agreement}).
A negative result is worded as what was screened shows: none among the
records screened, with the sample's interval for the rest.

\section{Conclusion}
\label{sec:reporting-conclusion}

The guidelines ask that limits on a search be stated and justified, that
a limited set of ranked results be noted with its rationale, that the
number of records per database and the dates be given, and that the
report say how many reviewers screened each record and what automation
did; the rapid-review guidance asks for a second screener on a part of
the records and an agreement check, and states that its proportion is
not evidence based.
None asks for the reason a screen stopped, which the procedure states
anyway.

